%% session.tex
%%
%% Copyright (C) 2000-2018 Simone Piccardi.  Permission is granted to
%% copy, distribute and/or modify this document under the terms of the GNU Free
%% Documentation License, Version 1.1 or any later version published by the
%% Free Software Foundation; with the Invariant Sections being "Un preambolo",
%% with no Front-Cover Texts, and with no Back-Cover Texts.  A copy of the
%% license is included in the section entitled "GNU Free Documentation
%% License".
%%

\chapter{Terminali e sessioni di lavoro}
\label{cha:session}


A lungo l'unico modo per interagire con sistema di tipo Unix è stato tramite
l'interfaccia dei terminali, ma anche oggi, nonostante la presenza di diverse
interfacce grafiche, essi continuano ad essere estensivamente usati per il
loro stretto legame la linea di comando.

Nella prima parte esamineremo i concetti base in cui si articola l'interfaccia
dei terminali, a partire dal sistema del \textit{job control} e delle sessioni
di lavoro, toccando infine anche le problematiche dell'interazione con
programmi non interattivi. Nella seconda parte tratteremo il funzionamento
dell'I/O su terminale, e delle varie peculiarità che esso viene ad assumere
nell'uso come interfaccia di accesso al sistema da parte degli utenti. La
terza parte coprirà le tematiche relative alla creazione e gestione dei
terminali virtuali, che consentono di replicare via software l'interfaccia dei
terminali.



\section{L'interazione con i terminali}
\label{sec:sess_job_control}

I terminali sono l'interfaccia con cui fin dalla loro nascita i sistemi
unix-like hanno gestito l'interazione con gli utenti, tramite quella riga di
comando che li caratterizza da sempre. Ma essi hanno anche una rilevanza
particolare perché quella dei terminali è l'unica interfaccia hardware usata
dal kernel per comunicare direttamente con gli utenti, con la cosiddetta
\textit{console} di sistema, senza dover passare per un programma.

Originariamente si trattava di dispositivi specifici (i terminali seriali, se
non addirittura le telescriventi). Oggi questa interfaccia viene in genere
emulata o tramite programmi o con le cosiddette console virtuali associate a
monitor e tastiera, ma esiste sempre la possibilità di associarla direttamente
ad alcuni dispositivi, come eventuali linee seriali, ed in certi casi, come
buona parte dei dispositivi embedded su cui gira Linux (come router, access
point, ecc.) questa resta anche l'unica opzione per una \textit{console} di
sistema.


\subsection{Il \textit{job control}}
\label{sec:sess_job_control_overview}

Viene comunemente chiamato \textit{job control} quell'insieme di funzionalità
il cui scopo è quello di permettere ad un utente di poter sfruttare le
capacità multitasking di un sistema Unix per eseguire in contemporanea più
processi, pur potendo accedere, di solito, ad un solo terminale, avendo cioè
un solo punto in cui si può avere accesso all'input ed all'output degli
stessi. Con le interfacce grafiche di \textit{X Window} e con i terminali
virtuali via rete oggi tutto questo non è più vero, dato che si può accedere a
molti terminali in contemporanea da una singola postazione di lavoro, ma il
sistema è nato prima dell'esistenza di tutto ciò.

Il \textit{job control} è una caratteristica opzionale, introdotta in BSD
negli anni '80, e successivamente standardizzata da POSIX.1. La sua
disponibilità nel sistema è verificabile attraverso il controllo della macro
\macro{\_POSIX\_JOB\_CONTROL}. In generale il \textit{job control} richiede il
supporto sia da parte della shell (quasi tutte ormai lo hanno), che da parte
del kernel. In particolare il kernel deve assicurare sia la presenza di un
driver per i terminali abilitato al \textit{job control} che quella dei
relativi segnali illustrati in sez.~\ref{sec:sig_job_control}. 

In un sistema che supporta il \textit{job control}, una volta completato il
login, l'utente avrà a disposizione una shell dalla quale eseguire i comandi e
potrà iniziare quella che viene chiamata una \textsl{sessione di lavoro}, che
riunisce (vedi sez.~\ref{sec:sess_proc_group}) tutti i processi eseguiti
all'interno dello stesso login (esamineremo tutto il processo in dettaglio in
sez.~\ref{sec:sess_login}).

Siccome la shell è collegata ad un solo terminale, che viene usualmente
chiamato \textsl{terminale di controllo}, (vedi sez.~\ref{sec:sess_ctrl_term})
un solo comando alla volta, quello che viene detto in \textit{foreground} o in
\textsl{primo piano}, potrà scrivere e leggere dal terminale. La shell però
può eseguire, aggiungendo una ``\cmd{\&}'' alla fine del comando, più
programmi in contemporanea, mandandoli come si dice, ``in
\textit{background}'' (letteralmente ``\textsl{sullo sfondo}''), nel qual caso
essi saranno eseguiti senza essere collegati al terminale.

Si noti come si sia parlato di comandi e non di programmi o processi. Fra le
funzionalità della shell infatti c'è anche quella di consentire di concatenare
più comandi in una sola riga con il \textit{pipelining}, ed in tal caso
verranno eseguiti più programmi. Inoltre, anche quando si invoca un singolo
programma, questo potrà sempre lanciare eventuali sotto-processi per eseguire
dei compiti specifici.

Per questo l'esecuzione di una riga di comando può originare più di un
processo, quindi nella gestione del \textit{job control} non si può far
riferimento ai singoli processi.  Per questo il kernel prevede la possibilità
di raggruppare più processi in un cosiddetto \textit{process group} (detto
anche \textsl{raggruppamento di processi}, vedi
sez.~\ref{sec:sess_proc_group}). Deve essere cura della shell far sì che tutti
i processi che originano da una stessa riga di comando appartengano allo
stesso raggruppamento di processi, in modo che le varie funzioni di controllo,
ed i segnali inviati dal terminale, possano fare riferimento ad esso.

In generale all'interno di una sessione avremo un eventuale (può non esserci)
\textit{process group} in \textit{foreground}, che riunisce i processi che
possono accedere al terminale, e più \textit{process group} in
\textit{background}, che non possono accedervi. Il \textit{job control}
prevede che quando un processo appartenente ad un raggruppamento in
\textit{background} cerca di accedere al terminale, venga inviato un segnale a
tutti i processi del raggruppamento, in modo da bloccarli (vedi
sez.~\ref{sec:sess_ctrl_term}).

Un comportamento analogo si ha anche per i segnali generati dai comandi di
tastiera inviati dal terminale, che vengono inviati a tutti i processi del
raggruppamento in \textit{foreground}. In particolare \cmd{C-z} interrompe
l'esecuzione del comando, che può poi essere mandato in \textit{background}
con il comando \cmd{bg}. Il comando \cmd{fg} consente invece di mettere in
\textit{foreground} un comando precedentemente lanciato in
\textit{background}. Si tenga presente che \cmd{bg} e \cmd{fg} sono comandi
interni alla shell, che non comportano l'esecuzione di un programma esterno,
ma operazioni di gestione compiute direttamente dalla shell stessa.

Di norma la shell si cura anche di notificare all'utente, di solito prima
della stampa a video del prompt, lo stato dei vari processi. Essa infatti sarà
in grado, grazie all'uso della funzione di sistema \func{waitpid} (vedi
sez.~\ref{sec:proc_wait}), di rilevare sia i processi che sono terminati, sia
i raggruppamenti che sono bloccati, in quest'ultimo caso si dovrà usare la
specifica opzione \const{WUNTRACED}, secondo quanto già illustrato in
sez.~\ref{sec:proc_wait}.


\subsection{I \textit{process group} e le \textsl{sessioni}}
\label{sec:sess_proc_group}

\itindbeg{process~group}

Come accennato in sez.~\ref{sec:sess_job_control_overview} nel job control i
processi vengono raggruppati in \textit{process group} e \textsl{sessioni};
per far questo vengono utilizzati due ulteriori identificatori (oltre quelli
visti in sez.~\ref{sec:proc_pid}) che il kernel associa a ciascun
processo:\footnote{in Linux questi identificatori sono mantenuti nei campi
  \var{pgrp} e \var{session} della struttura \kstruct{task\_struct} definita
  in \file{include/linux/sched.h}.}  l'identificatore del \textit{process
  group} e l'identificatore della \textsl{sessione}, che vengono indicati
rispettivamente con le sigle \ids{PGID} e \ids{SID}, e sono mantenuti in
variabili di tipo \type{pid\_t}. I valori di questi identificatori possono
essere visualizzati dal comando \cmd{ps} usando l'opzione \cmd{-j}.

Un \textit{process group} è pertanto definito da tutti i processi che hanno lo
stesso \ids{PGID}; è possibile leggere il valore di questo identificatore con
le funzioni di sistema \funcd{getpgid} e \funcd{getpgrp}, i cui prototipi
sono:

\begin{funcproto}{
\fhead{unistd.h}
\fdecl{pid\_t getpgid(pid\_t pid)} 
\fdesc{Legge il \ids{PGID} di un processo.} 
\fdecl{pid\_t getpgrp(void)}
\fdesc{Legge il \ids{PGID} del processo corrente.} 
}

{Le funzioni ritornano il \ids{PGID} richiesto in caso di successo,
  \func{getpgrp} ha sempre successo mentre \func{getpgid} restituisce $-1$ per
  un errore, nel qual caso \var{errno} potrà assumere solo il valore
  \errval{ESRCH} se il processo indicato non esiste.
}
\end{funcproto}

Le due funzioni sono definite nello standard POSIX.1-2001, ma la prima deriva
da SVr4 e la seconda da BSD4.2 dove però è previsto possa prendere un
argomento per indicare il \ids{PID} di un altro processo. Si può riottenere
questo comportamento se di definisce la macro \macro{\_BSD\_SOURCE} e non sono
definite le altre macro che richiedono la conformità a POSIX, X/Open o SysV
(vedi sez.~\ref{sec:intro_standard}).

La funzione \func{getpgid} permette di specificare il \ids{PID} del processo
di cui si vuole sapere il \ids{PGID}. Un valore nullo per \param{pid}
restituisce il \ids{PGID} del processo corrente, che è il comportamento
ordinario di \func{getpgrp}, che di norma equivalente a \code{getpgid(0)}.

In maniera analoga l'identificatore della sessione di un processo (il
\ids{SID}) può essere letto dalla funzione di sistema \funcd{getsid}, il cui
prototipo è:

\begin{funcproto}{
\fhead{unistd.h}
\fdecl{pid\_t getsid(pid\_t pid)}
\fdesc{Legge il \ids{SID} di un processo.} 
}

{La funzione ritorna l'identificatore (un numero positivo) in caso di successo
  e $-1$ per un errore, nel qual caso \var{errno} assumerà uno dei valori:
  \begin{errlist}
    \item[\errcode{EPERM}] il processo selezionato non fa parte della stessa
      sessione del processo corrente (solo in alcune implementazioni).
    \item[\errcode{ESRCH}] il processo selezionato non esiste.
  \end{errlist}
}
\end{funcproto}

La funzione è stata introdotta in Linux a partire dal kernel 1.3.44, il
supporto nelle librerie del C è iniziato dalla versione 5.2.19. La funzione
non era prevista originariamente da POSIX.1, che parla solo di processi leader
di sessione, e non di identificatori di sessione, ma è prevista da SVr4 e fa
parte di POSIX.1-2001.  Per poterla utilizzare occorre definire la macro
\macro{\_XOPEN\_SOURCE} ad un valore maggiore o uguale di 500. Su Linux
l'errore \errval{EPERM} non viene mai restituito.

Entrambi gli identificatori, \ids{SID} e \ids{PGID}, vengono inizializzati
nella creazione di ciascun processo con lo stesso valore che hanno nel
processo padre, per cui un processo appena creato appartiene sempre allo
stesso raggruppamento e alla stessa sessione del padre. Vedremo a breve come
sia possibile creare più \textit{process group} all'interno della stessa
sessione, e spostare i processi dall'uno all'altro, ma sempre all'interno di
una stessa sessione.

Ciascun raggruppamento di processi ha sempre un processo principale, il
cosiddetto \itindex{process~group~leader} \textit{process group leader} o più
brevemente \textit{group leader}, che è identificato dall'avere un \ids{PGID}
uguale al suo \ids{PID}. In genere questo è il primo processo del
raggruppamento, che si incarica di lanciare tutti gli altri. Un nuovo
raggruppamento si crea con la funzione di sistema \funcd{setpgrp}, il cui
prototipo è:

\begin{funcproto}{
\fhead{unistd.h}
\fdecl{int setpgrp(void)}
\fdesc{Rende un processo \textit{group leader} di un nuovo gruppo.} 
}

{La funzione ritorna il valore del nuovo \textit{process group} e non sono
  previsti errori.}
\end{funcproto}

La funzione assegna al \ids{PGID} il valore del \ids{PID} del processo
corrente, rendendolo in tal modo \textit{group leader} di un nuovo
raggruppamento. Tutti i successivi processi da esso creati apparterranno (a
meno di non cambiare di nuovo il \ids{PGID}) al nuovo raggruppamento.

La versione illustrata è quella usata nella definizione di POSIX.1, in BSD
viene usata una funzione con questo nome, che però è identica a
\func{setpgid}, che vedremo a breve, negli argomenti e negli effetti. Nella
\acr{glibc} viene sempre usata sempre questa definizione, a meno di non
richiedere esplicitamente la compatibilità all'indietro con BSD, definendo la
macro \macro{\_BSD\_SOURCE} ed evitando di definire le macro che richiedono
gli altri standard, come per \func{getpgrp}.

È inoltre possibile spostare un processo da un raggruppamento di processi ad
un altro cambiandone il \ids{PGID} con la funzione di sistema \funcd{setpgid},
il cui prototipo è:

\begin{funcproto}{
\fhead{unistd.h}
\fdecl{int setpgid(pid\_t pid, pid\_t pgid)}
\fdesc{Modifica il \ids{PGID} di un processo.} 
}

{La funzione ritorna il valore del nuovo \textit{process group}  in caso di
  successo e $-1$ per un errore, nel qual caso \var{errno} assumerà uno dei
  valori:
  \begin{errlist}
    \item[\errcode{EACCES}] il processo di cui si vuole cambiare il \ids{PGID}
      ha già eseguito una \func{exec}.
    \item[\errcode{EINVAL}] il valore di \param{pgid} è negativo.
    \item[\errcode{EPERM}] il cambiamento non è consentito.
    \item[\errcode{ESRCH}] il processo selezionato non esiste.
  \end{errlist}
}
\end{funcproto}


La funzione permette di cambiare il \ids{PGID} del processo indicato
dall'argomento \param{pid}, ma il cambiamento può essere effettuato solo se
l'argomento \param{pgid} indica un \textit{process group} che è nella stessa
sessione del processo chiamante.  Inoltre la funzione può essere usata
soltanto sul processo corrente o su uno dei suoi figli, ed in quest'ultimo
caso ha successo soltanto se questo non ha ancora eseguito una
\func{exec}.\footnote{questa caratteristica è implementata dal kernel che
  mantiene allo scopo un altro campo, \var{did\_exec}, nella struttura
  \kstruct{task\_struct}.}  Specificando un valore nullo per \param{pid} si
indica il processo corrente, mentre specificando un valore nullo
per \param{pgid} si imposta il \textit{process group} al valore del \ids{PID}
del processo selezionato, questo significa che \func{setpgrp} è equivalente a
\code{setpgid(0, 0)}.

Di norma questa funzione viene usata dalla shell quando si usano delle
pipeline, per mettere nello stesso \textit{process group} tutti i programmi
lanciati su ogni linea di comando; essa viene chiamata dopo una \func{fork}
sia dal processo padre, per impostare il valore nel figlio, che da
quest'ultimo, per sé stesso, in modo che il cambiamento di \textit{process
  group} sia immediato per entrambi; una delle due chiamate sarà ridondante,
ma non potendo determinare quale dei due processi viene eseguito per primo,
occorre eseguirle comunque entrambe per evitare di esporsi ad una \textit{race
  condition}.

Si noti come nessuna delle funzioni esaminate finora permetta di spostare un
processo da una sessione ad un altra; infatti l'unico modo di far cambiare
sessione ad un processo è quello di crearne una nuova con l'uso della funzione
di sistema \funcd{setsid}, il cui prototipo è:

\begin{funcproto}{
\fhead{unistd.h}
\fdecl{pid\_t setsid(void)}
\fdesc{Crea una nuova sessione sul processo corrente.} 
}

{La funzione ritorna il valore del nuovo \ids{SID} in caso di successo e $-1$
  per un errore, nel qual caso \var{errno} assumerà uno dei valori:
  \begin{errlist}
  \item[\errcode{EPERM}] il \ids{PGID} e \ids{PID} del processo coincidono.
  \end{errlist}
}
\end{funcproto}


La funzione imposta il \ids{PGID} ed il \ids{SID} del processo corrente al
valore del suo \ids{PID}, creando così una nuova sessione ed un nuovo
\textit{process group} di cui esso diventa leader (come per i \textit{process
  group} un processo si dice \textsl{leader di sessione} se il suo \ids{SID} è
uguale al suo \ids{PID}) ed unico componente.\footnote{in Linux la proprietà è
  mantenuta in maniera indipendente con un apposito campo \var{leader} in
  \kstruct{task\_struct}.} Inoltre la funzione distacca il processo da ogni
terminale di controllo (torneremo sull'argomento in
sez.~\ref{sec:sess_ctrl_term}) cui fosse in precedenza associato.

La funzione ha successo soltanto se il processo non è già leader di un
\textit{process group}, per cui per usarla di norma si esegue una \func{fork}
e si esce, per poi chiamare \func{setsid} nel processo figlio, in modo che,
avendo questo lo stesso \ids{PGID} del padre ma un \ids{PID} diverso, non ci
siano possibilità di errore.

Potrebbe sorgere il dubbio che, per il riutilizzo dei valori dei \ids{PID}
fatto nella creazione dei nuovi processi (vedi sez.~\ref{sec:proc_pid}), il
figlio venga ad assumere un valore corrispondente ad un \textit{process group}
esistente; questo viene evitato dal kernel che considera come disponibili per
un nuovo \ids{PID} solo valori che non corrispondono ad altri \ids{PID},
\ids{PGID} o \ids{SID} in uso nel sistema. Questa funzione viene usata di
solito nel processo di login (per i dettagli vedi sez.~\ref{sec:sess_login})
per raggruppare in una sessione tutti i comandi eseguiti da un utente dalla
sua shell.

\itindend{process~group}

\subsection{Il terminale di controllo e il controllo di sessione}
\label{sec:sess_ctrl_term}

Come accennato in sez.~\ref{sec:sess_job_control_overview}, nel sistema del
\textit{job control} i processi all'interno di una sessione fanno riferimento
ad un terminale di controllo (ad esempio quello su cui si è effettuato il
login), sul quale effettuano le operazioni di lettura e scrittura, e dal quale
ricevono gli eventuali segnali da tastiera. Nel caso di login grafico la cosa
può essere più complessa, e di norma l'I/O è effettuato tramite il server X,
ma ad esempio per i programmi, anche grafici, lanciati da un qualunque
emulatore di terminale, sarà quest'ultimo a fare da terminale (virtuale) di
controllo.

Per realizzare questa funzionalità lo standard POSIX.1 prevede che ad ogni
sessione possa essere associato un terminale di controllo; in Linux questo
viene realizzato mantenendo fra gli attributi di ciascun processo anche qual'è
il suo terminale di controllo.\footnote{lo standard POSIX.1 non specifica
  nulla riguardo l'implementazione; in Linux anch'esso viene mantenuto nella
  solita struttura \kstruct{task\_struct}, nel campo \var{tty}.}  In generale
ogni processo eredita dal padre, insieme al \ids{PGID} e al \ids{SID} anche il
terminale di controllo (vedi sez.~\ref{sec:proc_fork}). In questo modo tutti
processi originati dallo stesso leader di sessione mantengono lo stesso
terminale di controllo.

Alla creazione di una nuova sessione con \func{setsid} ogni associazione con
il precedente terminale di controllo viene cancellata, ed il processo che è
divenuto un nuovo leader di sessione dovrà essere associato ad un nuovo
terminale di controllo.\footnote{questo però solo quando necessario, cosa che,
  come vedremo in sez.~\ref{sec:sess_daemon}, non è sempre vera.} Questo viene
fatto dal kernel la prima volta che il processo apre un terminale (cioè uno
dei vari file di dispositivo \file{/dev/tty*}). In tal caso questo diventa
automaticamente il terminale di controllo, ed il processo diventa il
\textsl{processo di controllo} di quella sessione. Questo avviene
automaticamente a meno di non avere richiesto esplicitamente il contrario
aprendo il terminale di controllo con il flag \const{O\_NOCTTY} (vedi
sez.~\ref{sec:file_open_close}). In questo Linux segue la semantica di SVr4;
BSD invece richiede che il terminale venga allocato esplicitamente con una
\func{ioctl} con il comando \constd{TIOCSCTTY}.

In genere, a meno di redirezioni, nelle sessioni di lavoro il terminale di
controllo è associato ai file standard (input, output ed error) dei processi
nella sessione, ma solo quelli che fanno parte del cosiddetto raggruppamento
di \textit{foreground}, possono leggere e scrivere in certo istante. Per
impostare il raggruppamento di \textit{foreground} di un terminale si usa la
funzione \funcd{tcsetpgrp}, il cui prototipo è:

\begin{funcproto}{
\fhead{unistd.h}
\fhead{termios.h}
\fdecl{int tcsetpgrp(int fd, pid\_t pgrpid)}
\fdesc{Imposta il \textit{process group} di \textit{foreground}.}
}

{La funzione ritorna $0$ in caso di successo e $-1$ per un errore, nel qual
  caso \var{errno} assumerà uno dei valori: 
  \begin{errlist}
    \item[\errcode{ENOSYS}] il sistema non supporta il \textit{job control}.
    \item[\errcode{ENOTTY}] il file \param{fd} non corrisponde al terminale di
      controllo del processo chiamante.
    \item[\errcode{EPERM}] il \textit{process group} specificato non è nella
    stessa sessione del processo chiamante.
  \end{errlist}
  ed inoltre \errval{EBADF} ed \errval{EINVAL} nel loro significato generico.}
\end{funcproto}

La funzione imposta a \param{pgrpid} il \textit{process group} di
\textit{foreground} del terminale associato al file descriptor \param{fd}. La
funzione può essere eseguita con successo solo da un processo che
ha \param{fd} come terminale di controllo della propria sessione;
inoltre \param{pgrpid} deve essere un \textit{process group} (non vuoto)
appartenente alla stessa sessione del processo chiamante.

Come accennato in sez.~\ref{sec:sess_job_control_overview}, tutti i processi
(e relativi raggruppamenti) che non fanno parte del gruppo di
\textit{foreground} sono detti in \textit{background}; se uno si essi cerca di
accedere al terminale di controllo provocherà l'invio da parte del kernel di
uno dei due segnali \signal{SIGTTIN} o \signal{SIGTTOU} (a seconda che
l'accesso sia stato in lettura o scrittura) a tutto il suo \textit{process
  group}; dato che il comportamento di default di questi segnali (si riveda
quanto esposto in sez.~\ref{sec:sig_job_control}) è di fermare il processo, di
norma questo comporta che tutti i membri del gruppo verranno fermati, ma non
si avranno condizioni di errore.\footnote{la shell in genere notifica comunque
  un avvertimento, avvertendo la presenza di processi bloccati grazie all'uso
  di \func{waitpid}.} Se però si bloccano o ignorano i due segnali citati, le
funzioni di lettura e scrittura falliranno con un errore di \errcode{EIO}.

In genere la funzione viene chiamata da un processo che è gruppo di
\textit{foreground} per passare l'accesso al terminale ad altri processi,
quando viene chiamata da un processo che non è nel gruppo di
\textit{foreground}, a meno che questi non stia bloccando o ignorando il
segnale \signal{SIGTTOU}, detto segnale verrà inviato a tutti i processi del
gruppo di cui il chiamante fa parte.

Un processo può controllare qual è il gruppo di \textit{foreground} associato
ad un terminale con la funzione \funcd{tcgetpgrp}, il cui prototipo è:

\begin{funcproto}{
\fhead{unistd.h}
\fhead{termios.h}
\fdecl{pid\_t tcgetpgrp(int fd)}
\fdesc{Legge il \textit{process group} di \textit{foreground}.} 
}

{La funzione ritorna il \ids{PGID} del gruppo di \textit{foreground} in caso
  di successo e $-1$ per un errore, nel qual caso \var{errno} assumerà uno dei
  valori:
  \begin{errlist}
    \item[\errcode{ENOTTY}] non c'è un terminale di controllo o \param{fd} non
      corrisponde al terminale di controllo del processo chiamante.
  \end{errlist}
  ed inoltre \errval{EBADF} nel suo significato generico.}
\end{funcproto}

La funzione legge il \textit{process group} di \textit{foreground} associato
al file descriptor \param{fd}, che deve essere un terminale, restituendolo
come risultato. Sia questa funzione che la precedente sono state introdotte
con lo standard POSIX.1-2001, ma su Linux sono realizzate utilizzando le
operazioni di \func{ioctl} con i comandi \constd{TIOCGPGRP} e
\constd{TIOCSPGRP}.

Si noti come entrambe le funzioni usino come argomento il valore di un file
descriptor, il risultato comunque non dipende dal file descriptor che si usa
ma solo dal terminale cui fa riferimento. Il kernel inoltre permette a ciascun
processo di accedere direttamente al suo terminale di controllo attraverso il
file speciale \file{/dev/tty}, che per ogni processo è un sinonimo per il
proprio terminale di controllo. Questo consente anche a processi che possono
aver rediretto l'output di accedere al terminale di controllo, pur non
disponendo più del file descriptor originario; un caso tipico è il programma
\cmd{crypt} che accetta la redirezione sullo \textit{standard input} di un
file da decifrare, ma deve poi leggere la password dal terminale.

Un'altra caratteristica del terminale di controllo usata nel \textit{job
  control} è che utilizzando su di esso le combinazioni di tasti speciali
(\texttt{C-z}, \texttt{C-c}, \texttt{C-y} e \texttt{C-|}) si farà sì che il
kernel invii i corrispondenti segnali (rispettivamente \signal{SIGTSTP},
\signal{SIGINT}, \signal{SIGQUIT} e \signal{SIGTERM}, trattati in
sez.~\ref{sec:sig_job_control}) a tutti i processi del raggruppamento di
\textit{foreground}; in questo modo la shell può gestire il blocco e
l'interruzione dei vari comandi.

Per completare la trattazione delle caratteristiche del \textit{job control}
legate al terminale di controllo, occorre prendere in considerazione i vari
casi legati alla terminazione anomala dei processi, che sono di norma gestite
attraverso il segnale \signal{SIGHUP}. Il nome del segnale deriva da
\textit{hungup}, termine che viene usato per indicare la condizione in cui il
terminale diventa inutilizzabile, (letteralmente sarebbe
\textsl{impiccagione}).

Quando si verifica questa condizione, ad esempio se si interrompe la linea, o
va giù la rete o più semplicemente si chiude forzatamente la finestra di
terminale su cui si stava lavorando, il kernel provvederà ad inviare il
segnale di \signal{SIGHUP} al processo di controllo. L'azione preimpostata in
questo caso è la terminazione del processo, il problema che si pone è cosa
accade agli altri processi nella sessione, che non han più un processo di
controllo che possa gestire l'accesso al terminale, che potrebbe essere
riutilizzato per qualche altra sessione.

Lo standard POSIX.1 prevede che quando il processo di controllo termina, che
ciò avvenga o meno per un \textit{hungup} del terminale (ad esempio si
potrebbe terminare direttamente la shell con \cmd{kill}) venga inviato un
segnale di \signal{SIGHUP} ai processi del raggruppamento di
\textit{foreground}. In questo modo essi potranno essere avvisati che non
esiste più un processo in grado di gestire il terminale (di norma tutto ciò
comporta la terminazione anche di questi ultimi).

Restano però gli eventuali processi in \textit{background}, che non ricevono
il segnale; in effetti se il terminale non dovesse più servire essi potrebbero
proseguire fino al completamento della loro esecuzione; ma si pone il problema
di come gestire quelli che sono bloccati, o che si bloccano nell'accesso al
terminale, in assenza di un processo che sia in grado di effettuare il
controllo dello stesso.

Questa è la situazione in cui si ha quello che viene chiamato un
\itindex{process~group~orphaned} \textit{orphaned process group}. Lo standard
POSIX.1 lo definisce come un \textit{process group} i cui processi hanno come
padri esclusivamente o altri processi nel raggruppamento, o processi fuori
della sessione.  Lo standard prevede inoltre che se la terminazione di un
processo fa sì che un raggruppamento di processi diventi orfano e se i suoi
membri sono bloccati, ad essi vengano inviati in sequenza i segnali di
\signal{SIGHUP} e \signal{SIGCONT}.

La definizione può sembrare complicata, e a prima vista non è chiaro cosa
tutto ciò abbia a che fare con il problema della terminazione del processo di
controllo.  Consideriamo allora cosa avviene di norma nel \textit{job
  control}: una sessione viene creata con \func{setsid} che crea anche un
nuovo \textit{process group}. Per definizione quest'ultimo è sempre
\textsl{orfano}, dato che il padre del leader di sessione è fuori dalla stessa
e il nuovo \textit{process group} contiene solo il leader di sessione. Questo
è un caso limite, e non viene emesso nessun segnale perché quanto previsto
dallo standard riguarda solo i raggruppamenti che diventano orfani in seguito
alla terminazione di un processo.\footnote{l'emissione dei segnali infatti
  avviene solo nella fase di uscita del processo, come una delle operazioni
  legate all'esecuzione di \func{\_exit}, secondo quanto illustrato in
  sez.~\ref{sec:proc_termination}.}

Il leader di sessione provvederà a creare nuovi raggruppamenti che a questo
punto non sono orfani in quanto esso resta padre per almeno uno dei processi
del gruppo (gli altri possono derivare dal primo). Alla terminazione del
leader di sessione però avremo che, come visto in
sez.~\ref{sec:proc_termination}, tutti i suoi figli vengono adottati da
\cmd{init}, che è fuori dalla sessione. Questo renderà orfani tutti i
\textit{process group} creati direttamente dal leader di sessione a meno di
non aver spostato con \func{setpgid} un processo da un gruppo ad un altro,
(cosa che di norma non viene fatta) i quali riceveranno, nel caso siano
bloccati, i due segnali; \signal{SIGCONT} ne farà proseguire l'esecuzione, ed
essendo stato nel frattempo inviato anche \signal{SIGHUP}, se non c'è un
gestore per quest'ultimo, i processi bloccati verranno automaticamente
terminati.



\subsection{Dal login alla shell}
\label{sec:sess_login}

L'organizzazione del sistema del job control è strettamente connessa alle
modalità con cui un utente accede al sistema per dare comandi, collegandosi ad
esso con un terminale, che sia questo realmente tale, come un VT100 collegato
ad una seriale o virtuale, come quelli associati a schermo e tastiera o ad una
connessione di rete. Dato che i concetti base sono gli stessi, e dato che alla
fine le differenze sono nel dispositivo cui il kernel associa i file standard
(vedi tab.~\ref{tab:file_std_files}) per l'I/O, tratteremo solo il caso
classico del terminale, in generale nel caso di login via rete o di terminali
lanciati dall'interfaccia grafica cambia anche il processo da cui ha origine
l'esecuzione della shell.

Abbiamo già brevemente illustrato in sez.~\ref{sec:intro_kern_and_sys} le
modalità con cui il sistema si avvia, e di come, a partire da \cmd{init},
vengano lanciati tutti gli altri processi. Adesso vedremo in maniera più
dettagliata le modalità con cui il sistema arriva a fornire ad un utente la
shell che gli permette di lanciare i suoi comandi su un terminale.

Nella maggior parte delle distribuzioni di GNU/Linux viene usata la procedura
di avvio di System V;\footnote{in realtà negli ultimi tempi questa situazione
  sta cambiando, e sono state proposte diversi possibili rimpiazzi per il
  tradizionale \texttt{init} di System V, come \texttt{upstart} o
  \texttt{systemd}, ma per quanto trattato in questa sezione il risultato
  finale non cambia, si avrà comunque il lancio di un programma che consenta
  l'accesso al terminale.}  questa prevede che \cmd{init} legga dal file di
configurazione \conffiled{/etc/inittab} quali programmi devono essere lanciati,
ed in quali modalità, a seconda del cosiddetto \textit{run level}, anch'esso
definito nello stesso file.

Tralasciando la descrizione del sistema dei \textit{run level}, (per il quale
si rimanda alla lettura delle pagine di manuale di \cmd{init} e di
\file{inittab} o alla trattazione in sez.~5.3 di \cite{AGL}) quello che
comunque viene sempre fatto è di eseguire almeno una istanza di un programma
che permetta l'accesso ad un terminale. Uno schema di massima della procedura
è riportato in fig.~\ref{fig:sess_term_login}.

\begin{figure}[!htb]
  \centering \includegraphics[width=13cm]{img/tty_login}
  \caption{Schema della procedura di login su un terminale.}
  \label{fig:sess_term_login}
\end{figure}

Un terminale, che esso sia un terminale effettivo, attaccato ad una seriale o
ad un altro tipo di porta di comunicazione, o una delle console virtuali
associate allo schermo, o un terminale virtuale ad uso degli emulatori o delle
sessioni di rete, viene sempre visto attraverso un apposito file di
dispositivo che presenta una serie di caratteristiche comuni che vanno a
costituire l'interfaccia generica di accesso ai terminali.

Per controllare un terminale fisico come la seriale o le console virtuali
dello schermo si usa di solito il programma \cmd{getty} (o una delle sue
varianti). Alla radice della catena che porta ad una shell per i comandi
perciò c'è sempre \cmd{init} che esegue prima una \func{fork} e poi una
\func{exec} per lanciare una istanza di questo programma, il tutto ripetuto
per ciascuno dei terminali che si vogliono attivare, secondo quanto indicato
dall'amministratore nel file di configurazione del programma,
\conffile{/etc/inittab}.

Quando viene lanciato da \cmd{init} il programma parte con i privilegi di
amministratore e con un ambiente vuoto; \cmd{getty} si cura di chiamare
\func{setsid} per creare una nuova sessione ed un nuovo \textit{process
  group}, e di aprire il terminale (che così diventa il terminale di controllo
della sessione) in lettura sullo \textit{standard input} ed in scrittura sullo
\textit{standard output} e sullo \textit{standard error}; inoltre effettuerà,
qualora servano, ulteriori impostazioni.\footnote{ad esempio, come qualcuno si
  sarà accorto scrivendo un nome di login in maiuscolo, può effettuare la
  conversione automatica dell'input in minuscolo, ponendosi in una modalità
  speciale che non distingue fra i due tipi di caratteri (a beneficio di
  alcuni vecchi terminali che non supportavano le minuscole).}  Alla fine il
programma stamperà un messaggio di benvenuto per poi porsi in attesa
dell'immissione del nome di un utente.

Una volta che si sia immesso il nome di login \cmd{getty} esegue direttamente
il programma \cmd{login} con una \func{execle}, passando come argomento la
stringa con il nome, ed un ambiente opportunamente costruito che contenga
quanto necessario; ad esempio di solito viene opportunamente inizializzata la
variabile di ambiente \envvar{TERM} per identificare il terminale su cui si
sta operando, a beneficio dei programmi che verranno lanciati in seguito.

A sua volta \cmd{login}, che mantiene i privilegi di amministratore, usa il
nome dell'utente per effettuare una ricerca nel database degli utenti (in
genere viene chiamata \func{getpwnam}, che abbiamo visto in
sez.~\ref{sec:sys_user_group}, per leggere la password e gli altri dati dal
database degli utenti) e richiede una password. Se l'utente non esiste o se la
password non corrisponde\footnote{il confronto non viene effettuato con un
  valore in chiaro; quanto immesso da terminale viene invece a sua volta
  criptato, ed è il risultato che viene confrontato con il valore che viene
  mantenuto nel database degli utenti.} la richiesta viene ripetuta un certo
numero di volte dopo di che \cmd{login} esce ed \cmd{init} provvede a
rilanciare un'altra istanza di \cmd{getty}.

Se invece la password corrisponde \cmd{login} esegue \func{chdir} per
impostare come directory di lavoro la \textit{home directory} dell'utente,
cambia i diritti di accesso al terminale (con \func{chown} e \func{chmod}) per
assegnarne la titolarità all'utente ed al suo gruppo principale, assegnandogli
al contempo i diritti di lettura e scrittura.\footnote{oggi queste operazioni,
  insieme ad altre relative alla contabilità ed alla tracciatura degli
  accessi, vengono gestite dalle distribuzioni più recenti in una maniera
  generica appoggiandosi a servizi di sistema come \textit{ConsoleKit}, ma il
  concetto generale resta sostanzialmente lo stesso.}  Inoltre il programma
provvede a costruire gli opportuni valori per le variabili di ambiente, come
\envvar{HOME}, \envvar{SHELL}, ecc.  Infine attraverso l'uso di \func{setuid},
\func{setgid} e \func{initgroups} verrà cambiata l'identità del proprietario
del processo, infatti, come spiegato in sez.~\ref{sec:proc_setuid}, avendo
invocato tali funzioni con i privilegi di amministratore, tutti gli \ids{UID}
ed i \ids{GID} (reali, effettivi e salvati) saranno impostati a quelli
dell'utente.

A questo punto \cmd{login} provvederà (fatte salve eventuali altre azioni
iniziali, come la stampa di messaggi di benvenuto o il controllo della posta)
ad eseguire con un'altra \func{exec} la shell, che si troverà con un ambiente
già pronto con i file standard di tab.~\ref{tab:file_std_files} impostati sul
terminale, e pronta, nel ruolo di leader di sessione e di processo di
controllo per il terminale, a gestire l'esecuzione dei comandi come illustrato
in sez.~\ref{sec:sess_job_control_overview}.

Dato che il processo padre resta sempre \cmd{init} quest'ultimo potrà
provvedere, ricevendo un \signal{SIGCHLD} all'uscita della shell quando la
sessione di lavoro è terminata, a rilanciare \cmd{getty} sul terminale per
ripetere da capo tutto il procedimento.



\subsection{Interazione senza terminale: i \textsl{demoni} ed il
  \textit{syslog}}
\label{sec:sess_daemon}

Come sottolineato fin da sez.~\ref{sec:intro_base_concept}, in un sistema
unix-like tutte le operazioni sono eseguite tramite processi, comprese quelle
operazioni di sistema (come l'esecuzione dei comandi periodici, o la consegna
della posta, ed in generale tutti i programmi di servizio) che non hanno
niente a che fare con la gestione diretta dei comandi dell'utente.

Questi programmi, che devono essere eseguiti in modalità non interattiva e
senza nessun intervento dell'utente, sono normalmente chiamati
\textsl{demoni}, (o \textit{daemons}), nome ispirato dagli omonimi spiritelli
della mitologia greca che svolgevano compiti che gli dei trovavano noiosi, di
cui parla anche Socrate (che sosteneva di averne uno al suo servizio).

%TODO ricontrollare, i miei ricordi di filosofia sono piuttosto datati.

Se però si lancia un programma demone dalla riga di comando in un sistema che
supporta, come Linux, il \textit{job control} esso verrà comunque associato ad
un terminale di controllo e mantenuto all'interno di una sessione, e anche se
può essere mandato in background e non eseguire più nessun I/O su terminale,
si avranno comunque tutte le conseguenze che abbiamo trattato in
sez.~\ref{sec:sess_ctrl_term}, in particolare l'invio dei segnali in
corrispondenza dell'uscita del leader di sessione.

Per questo motivo un programma che deve funzionare come demone deve sempre
prendere autonomamente i provvedimenti opportuni (come distaccarsi dal
terminale e dalla sessione) ad impedire eventuali interferenze da parte del
sistema del \textit{job control}; questi sono riassunti in una lista di
prescrizioni\footnote{ad esempio sia Stevens in \cite{APUE}, che la
  \textit{Unix Programming FAQ} \cite{UnixFAQ} ne riportano di sostanzialmente
  identiche.} da seguire quando si scrive un demone.

Pertanto, quando si lancia un programma che deve essere eseguito come demone
occorrerà predisporlo in modo che esso compia le seguenti azioni:
\begin{enumerate*}
\item Eseguire una \func{fork} e terminare immediatamente il processo padre
  proseguendo l'esecuzione nel figlio.  In questo modo si ha la certezza che
  il figlio non è un \textit{process group leader}, (avrà il \ids{PGID} del
  padre, ma un \ids{PID} diverso) e si può chiamare \func{setsid} con
  successo. Inoltre la shell considererà terminato il comando all'uscita del
  padre.
\item Eseguire \func{setsid} per creare una nuova sessione ed un nuovo
  raggruppamento di cui il processo diventa automaticamente il leader, che
  però non ha associato nessun terminale di controllo.
\item Assicurarsi che al processo non venga associato in seguito nessun nuovo
  terminale di controllo; questo può essere fatto sia avendo cura di usare
  sempre l'opzione \const{O\_NOCTTY} nell'aprire i file di terminale, che
  eseguendo una ulteriore \func{fork} uscendo nel padre e proseguendo nel
  figlio. In questo caso, non essendo più quest'ultimo un leader di sessione
  non potrà ottenere automaticamente un terminale di controllo.
\item Eseguire una \func{chdir} per impostare la directory di lavoro del
  processo (su \file{/} o su una directory che contenga dei file necessari per
  il programma), per evitare che la directory da cui si è lanciato il processo
  resti in uso e non sia possibile rimuoverla o smontare il filesystem che la
  contiene.
\item Impostare la maschera dei permessi (di solito con \code{umask(0)}) in
  modo da non essere dipendenti dal valore ereditato da chi ha lanciato
  originariamente il processo.
\item Chiudere tutti i file aperti che non servono più (in generale tutti); in
  particolare vanno chiusi i file standard che di norma sono ancora associati
  al terminale (un'altra opzione è quella di redirigerli verso
  \file{/dev/null}).
\end{enumerate*}


In Linux buona parte di queste azioni possono venire eseguite invocando la
funzione \funcd{daemon} (fornita dalla \acr{glibc}), introdotta per la prima
volta in BSD4.4; il suo prototipo è:

\begin{funcproto}{
\fhead{unistd.h}
\fdecl{int daemon(int nochdir, int noclose)}
\fdesc{Rende il processo un demone.} 
}

{La funzione ritorna (nel nuovo processo) $0$ in caso di successo e $-1$ per
  un errore, nel qual caso \var{errno} assumerà uno dei valori i
    valori impostati dalle sottostanti \func{fork} e \func{setsid}.}
\end{funcproto}

La funzione esegue una \func{fork}, per uscire subito, con \func{\_exit}, nel
padre, mentre l'esecuzione prosegue nel figlio che esegue subito una
\func{setsid}. In questo modo si compiono automaticamente i passi 1 e 2 della
precedente lista. Se \param{nochdir} è nullo la funzione imposta anche la
directory di lavoro su \file{/}, se \param{noclose} è nullo i file standard
vengono rediretti su \file{/dev/null} (corrispondenti ai passi 4 e 6); in caso
di valori non nulli non viene eseguita nessuna altra azione.

Dato che un programma demone non può più accedere al terminale, si pone il
problema di come fare per la notifica di eventuali errori, non potendosi più
utilizzare lo \textit{standard error}. Per il normale I/O infatti ciascun
demone avrà le sue modalità di interazione col sistema e gli utenti a seconda
dei compiti e delle funzionalità che sono previste; ma gli errori devono
normalmente essere notificati all'amministratore del sistema.

\itindbeg{syslog}

Una soluzione può essere quella di scrivere gli eventuali messaggi su uno
specifico file (cosa che a volte viene fatta comunque) ma questo comporta il
grande svantaggio che l'amministratore dovrà tenere sotto controllo un file
diverso per ciascun demone, e che possono anche generarsi conflitti di nomi.
Per questo in BSD4.2 venne introdotto un servizio di sistema, il
\textit{syslog}, che oggi si trova su tutti i sistemi Unix, e che permette ai
demoni di inviare messaggi all'amministratore in una maniera standardizzata.

Il servizio prevede vari meccanismi di notifica, e, come ogni altro servizio
in un sistema unix-like, viene gestito attraverso un apposito programma, che è
anch'esso un \textsl{demone}. In generale i messaggi di errore vengono
raccolti dal file speciale \file{/dev/log}, un socket locale (vedi
sez.~\ref{sec:sock_sa_local}) dedicato a questo scopo, o via rete, con un
socket UDP e trattati dal demone che gestisce il servizio. Il più comune di
questi è \texttt{syslogd}, che consente un semplice smistamento dei messaggi
sui file in base alle informazioni in esse presenti; oggi però
\texttt{syslogd} è in sostanziale disuso, sostituito da programmi più
sofisticati come \texttt{rsyslog} o \texttt{syslog-ng}.

Il servizio del \textit{syslog} permette infatti di trattare i vari messaggi
classificandoli attraverso due indici: il primo, chiamato \textit{facility},
suddivide in diverse categorie i messaggi in modo di raggruppare quelli
provenienti da operazioni che hanno attinenza fra loro, ed è organizzato in
sottosistemi (kernel, posta elettronica, demoni di stampa, ecc.). Il secondo,
chiamato \textit{priority}, identifica l'importanza dei vari messaggi, e
permette di classificarli e differenziare le modalità di notifica degli
stessi.

Il sistema del \textit{syslog} attraverso il proprio demone di gestione
provvede poi a riportare i messaggi all'amministratore attraverso una serie
differenti meccanismi come:
\begin{itemize*}
\item scriverli su un file (comunemente detto \textit{log file}, o giornale),
\item scriverli sulla console,
\item scriverli sui terminali degli utenti connessi,
\item inviarli via mail ad uno specifico utente,
\item inviarli ad un altro programma,
\item inviarli via rete ad una macchina di raccolta,
\item ignorarli completamente;
\end{itemize*}
le modalità con cui queste azioni vengono realizzate dipendono ovviamente dal
demone che si usa, per la gestione del quale si rimanda ad un testo di
amministrazione di sistema.\footnote{l'argomento è ad esempio coperto dal
  capitolo 3.2.3 si \cite{AGL}.}

La \acr{glibc} definisce una serie di funzioni standard con cui un processo
può accedere in maniera generica al servizio di \textit{syslog}, che però
funzionano solo localmente; se si vogliono inviare i messaggi ad un altro
sistema occorre farlo esplicitamente con un socket UDP, o utilizzare le
capacità di reinvio del servizio.

La prima funzione definita dall'interfaccia è \funcd{openlog}, che inizializza
una connessione al servizio di \textit{syslog}. Essa in generale non è
necessaria per l'uso del servizio, ma permette di impostare alcuni valori che
controllano gli effetti delle chiamate successive; il suo prototipo è:

\begin{funcproto}{
\fhead{syslog.h}
\fdecl{void openlog(const char *ident, int option, int facility)}
\fdesc{Inizializza una connessione al sistema del \textit{syslog}.} 
}

{La funzione non restituisce nulla.}
\end{funcproto}

La funzione permette di specificare, tramite \param{ident}, l'identità di chi
ha inviato il messaggio (di norma si passa il nome del programma, come
specificato da \code{argv[0]}), e la stringa verrà preposta all'inizio di ogni
messaggio. Si tenga presente che il valore di \param{ident} che si passa alla
funzione è un puntatore, se la stringa cui punta viene cambiata lo sarà pure
nei successivi messaggi, e se viene cancellata i risultati potranno essere
impredicibili, per questo è sempre opportuno usare una stringa costante.

L'argomento \param{facility} permette invece di preimpostare per le successive
chiamate l'omonimo indice che classifica la categoria del messaggio.
L'argomento è interpretato come una maschera binaria, e pertanto è possibile
inviare i messaggi su più categorie alla volta. I valori delle costanti che
identificano ciascuna categoria sono riportati in
tab.~\ref{tab:sess_syslog_facility}, il valore di \param{facility} deve essere
specificato con un OR aritmetico.

\begin{table}[htb]
  \footnotesize
  \centering
  \begin{tabular}[c]{|l|p{8cm}|}
    \hline
    \textbf{Valore}& \textbf{Significato}\\
    \hline
    \hline
    \constd{LOG\_AUTH}     & Messaggi relativi ad autenticazione e sicurezza,
                             obsoleto, è sostituito da \const{LOG\_AUTHPRIV}.\\
    \constd{LOG\_AUTHPRIV} & Sostituisce \const{LOG\_AUTH}.\\
    \constd{LOG\_CRON}     & Messaggi dei demoni di gestione dei comandi
                             programmati (\cmd{cron} e \cmd{at}).\\
    \constd{LOG\_DAEMON}   & Demoni di sistema.\\
    \constd{LOG\_FTP}      & Servizio FTP.\\
    \constd{LOG\_KERN}     & Messaggi del kernel.\\
    \constd{LOG\_LOCAL0}   & Riservato all'amministratore per uso locale.\\
    \hspace{.5cm}$\vdots$ &   \hspace{3cm}$\vdots$\\
    \constd{LOG\_LOCAL7}   & Riservato all'amministratore per uso locale.\\
    \constd{LOG\_LPR}      & Messaggi del sistema di gestione delle stampanti.\\
    \constd{LOG\_MAIL}     & Messaggi del sistema di posta elettronica.\\
    \constd{LOG\_NEWS}     & Messaggi del sistema di gestione delle news 
                            (USENET).\\
    \constd{LOG\_SYSLOG}   & Messaggi generati dal demone di gestione del
                            \textit{syslog}.\\
    \constd{LOG\_USER}     & Messaggi generici a livello utente.\\
    \constd{LOG\_UUCP}     & Messaggi del sistema UUCP (\textit{Unix to Unix
                             CoPy}), ormai in disuso.\\
\hline
\end{tabular}
\caption{Valori possibili per l'argomento \param{facility} di \func{openlog}.}
\label{tab:sess_syslog_facility}
\end{table}

L'argomento \param{option} serve invece per controllare il comportamento della
funzione \func{openlog} e delle modalità con cui le successive chiamate
scriveranno i messaggi, esso viene specificato come maschera binaria composta
con un OR aritmetico di una qualunque delle costanti riportate in
tab.~\ref{tab:sess_openlog_option}.

\begin{table}[htb]
  \footnotesize
\centering
\begin{tabular}[c]{|l|p{8cm}|}
\hline
\textbf{Valore}& \textbf{Significato}\\
\hline
\hline
  \constd{LOG\_CONS}   & Scrive sulla console in caso di errore nell'invio del
                         messaggio al sistema del \textit{syslog}. \\
  \constd{LOG\_NDELAY} & Apre la connessione al sistema del \textit{syslog}
                         subito invece di attendere l'invio del primo
                         messaggio.\\ 
  \constd{LOG\_NOWAIT} & Non usato su Linux, su altre piattaforme non attende i
                         processi figli creati per inviare il messaggio.\\
  \constd{LOG\_ODELAY} & Attende il primo messaggio per aprire la connessione al
                         sistema del \textit{syslog}.\\ 
  \constd{LOG\_PERROR} & Stampa anche su \file{stderr} (non previsto in
                         POSIX.1-2001).\\ 
  \constd{LOG\_PID}    & Inserisce nei messaggi il \ids{PID} del processo
                         chiamante.\\
\hline
\end{tabular}
\caption{Valori possibili per l'argomento \param{option} di \func{openlog}.}
\label{tab:sess_openlog_option}
\end{table}

La funzione che si usa per generare un messaggio è \funcd{syslog}, dato che
l'uso di \func{openlog} è opzionale, sarà quest'ultima a provvede a chiamare la
prima qualora ciò non sia stato fatto (nel qual caso il valore di
\param{ident} è \val{NULL}). Il suo prototipo è:

\begin{funcproto}{
\fhead{syslog.h}
\fdecl{void syslog(int priority, const char *format, ...)}
\fdesc{Genera un messaggio per il \textit{syslog}.} 
}

{La funzione non restituisce nulla.}
\end{funcproto}

La funzione genera un messaggio le cui caratteristiche sono indicate
da \param{priority}. Per i restanti argomenti il suo comportamento è analogo a
quello di \func{printf}, e il valore dell'argomento \param{format} è identico
a quello descritto nella pagina di manuale di quest'ultima (per i valori
principali si può vedere la trattazione sommaria che se ne è fatto in
sez.~\ref{sec:file_formatted_io}).  L'unica differenza è che la sequenza
\val{\%m} viene rimpiazzata dalla stringa restituita da
\code{strerror(errno)}. Gli argomenti seguenti i primi due devono essere
forniti secondo quanto richiesto da \param{format}.

\begin{table}[htb]
  \footnotesize
  \centering
  \begin{tabular}[c]{|l|p{8cm}|}
    \hline
    \textbf{Valore}& \textbf{Significato}\\
    \hline
    \hline
    \constd{LOG\_EMERG}   & Il sistema è inutilizzabile.\\
    \constd{LOG\_ALERT}   & C'è una emergenza che richiede intervento
                            immediato.\\
    \constd{LOG\_CRIT}    & Si è in una condizione critica.\\
    \constd{LOG\_ERR}     & Si è in una condizione di errore.\\
    \constd{LOG\_WARNING} & Messaggio di avvertimento.\\
    \constd{LOG\_NOTICE}  & Notizia significativa relativa al comportamento.\\
    \constd{LOG\_INFO}    & Messaggio informativo.\\
    \constd{LOG\_DEBUG}   & Messaggio di debug.\\
    \hline
  \end{tabular}
  \caption{Valori possibili per l'indice di importanza del messaggio da
    specificare nell'argomento \param{priority} di \func{syslog}.}
  \label{tab:sess_syslog_priority}
\end{table}

L'argomento \param{priority} permette di impostare sia la \textit{facility}
che la \textit{priority} del messaggio. In realtà viene prevalentemente usato
per specificare solo quest'ultima in quanto la prima viene di norma
preimpostata con \func{openlog}. La priorità è indicata con un valore numerico
specificabile attraverso le costanti riportate in
tab.~\ref{tab:sess_syslog_priority}.  

La \acr{glibc}, seguendo POSIX.1-2001, prevede otto diverse priorità
ordinate da 0 a 7, in ordine di importanza decrescente; questo comporta che i
tre bit meno significativi dell'argomento \param{priority} sono occupati da
questo valore, mentre i restanti bit più significativi vengono usati per
specificare la \textit{facility}.  Nel caso si voglia specificare anche la
\textit{facility} basta eseguire un OR aritmetico del valore della priorità
con la maschera binaria delle costanti di tab.~\ref{tab:sess_syslog_facility}.

Una funzione sostanzialmente identica a \func{syslog} è \funcd{vsyslog}. La
funzione è originaria di BSD e per utilizzarla deve essere definito
\macro{\_BSD\_SOURCE}; il suo prototipo è:

\begin{funcproto}{
\fhead{syslog.h}
\fdecl{void vsyslog(int priority, const char *format, va\_list src)}
\fdesc{Genera un messaggio per il \textit{syslog}.} 
}

{La funzione non restituisce nulla.}
\end{funcproto}

La sola differenza con \func{syslog} è quella di prendere invece di una lista
di argomenti esplicita un unico argomento finale passato nella forma di una
\macro{va\_list}; la funzione risulta utile qualora si ottengano gli argomenti
dalla invocazione di un'altra funzione \textit{variadic} (si ricordi quanto
visto in sez.~\ref{sec:proc_variadic}).

Per semplificare la gestione della scelta del livello di priorità a partire
dal quale si vogliono registrare i messaggi, le funzioni di gestione
mantengono per ogni processo una maschera che determina quale delle chiamate
effettuate a \func{syslog} verrà effettivamente registrata.  In questo modo
sarà possibile escludere i livelli di priorità che non interessa registrare,
impostando opportunamente la maschera una volta per tutte.

Questo significa che in genere nei programmi vengono comunque previste le
chiamate a \func{syslog} per tutti i livelli di priorità, ma poi si imposta
questa maschera per registrare solo quello che effettivamente interessa. La
funzione che consente di fare questo è \funcd{setlogmask}, ed il suo prototipo
è:

\begin{funcproto}{
\fhead{syslog.h}
\fdecl{int setlogmask(int mask)}
\fdesc{Imposta la maschera dei messaggi del \textit{syslog}.} 
}

{La funzione ritorna il precedente valore della maschera dei messaggi e non
  prevede errori.}
\end{funcproto}

La funzione restituisce il valore della maschera corrente, e se si passa un
valore nullo per \param{mask} la maschera corrente non viene modificata; in
questo modo si può leggere il valore della maschera corrente. Indicando un
valore non nullo per \param{mask} la registrazione dei messaggi viene
disabilitata per tutte quelle priorità che non rientrano nella maschera.

In genere il valore viene impostato usando la macro
\macrod{LOG\_MASK}\texttt{(p)} dove \code{p} è una delle costanti di
tab.~\ref{tab:sess_syslog_priority}. É inoltre disponibile anche la macro
\macrod{LOG\_UPTO}\texttt{(p)} che permette di specificare automaticamente
tutte le priorità fino a quella indicata da \code{p}.

Una volta che si sia certi che non si intende registrare più nessun messaggio
si può chiudere esplicitamente la connessione al \textit{syslog} (l'uso di
questa funzione è comunque completamente opzionale) con la funzione
\funcd{closelog}, il cui prototipo è:

\begin{funcproto}{
\fhead{syslog.h}
\fdecl{void closelog(void)}
\fdesc{Chiude la connessione al \textit{syslog}.} 
}

{La funzione non ritorna nulla.}
\end{funcproto}


Come si evince anche dalla presenza della facility \const{LOG\_KERN} in
tab.~\ref{tab:sess_syslog_facility}, uno dei possibili utenti del servizio del
\textit{syslog} è anche il kernel, che a sua volta può avere necessità di
inviare messaggi verso l'\textit{user space}. I messaggi del kernel sono
mantenuti in un apposito buffer circolare e generati all'interno del kernel
tramite la funzione \texttt{printk}, analoga alla \func{printf} usata in
\textit{user space}.\footnote{una trattazione eccellente dell'argomento si
  trova nel quarto capitolo di \cite{LinDevDri}.}

Come per i messaggi ordinari anche i messaggi del kernel hanno una priorità ma
in questo caso non si può contare sulla coincidenza con le costanti di
tab.~\ref{tab:sess_syslog_priority} dato che il codice del kernel viene
mantenuto in maniera indipendente dalle librerie del C. Per questo motivo le
varie priorità usate dal kernel sono associate ad un valore numerico che viene
tradotto in una stringa preposta ad ogni messaggio, secondo i valori che si
sono riportati in fig.~\ref{fig:printk_priority}

\begin{figure}[!htb]
  \footnotesize \centering
  \begin{minipage}[c]{0.80\textwidth}
    \includestruct{listati/printk_prio.c}
  \end{minipage} 
  \normalsize 
  \caption{Definizione delle stringhe coi relativi valori numerici che
    indicano le priorità dei messaggi del kernel (ripresa da
    \file{include/linux/kernel.h}).}
  \label{fig:printk_priority}
\end{figure}

Dato che i messaggi generati da \texttt{printk} hanno un loro specifico
formato tradizionalmente si usava un demone ausiliario, \cmd{klogd}, per
leggerli, rimappare le loro priorità sui valori di
tab.~\ref{tab:sess_syslog_priority} ed inviarli al sistema del \textit{syslog}
nella facility \const{LOG\_KERN}.  Oggi i nuovi demoni più avanzati che
realizzano il servizio (come \texttt{rsyslog} o \texttt{syslog-ng}) sono in
grado di fare tutto questo da soli leggendoli direttamente senza necessità di
un intermediario.

Ma i messaggi del kernel non sono necessariamente connessi al sistema del
\textit{syslog}; ad esempio possono anche essere letti direttamente dal buffer
circolare con il comando \texttt{dmesg}. Inoltre è previsto che se superano
una certa priorità essi vengano stampati direttamente sul terminale indicato
come \textit{console} di sistema,\footnote{quello che viene indicato con il
  parametro di avvio \texttt{console} del kernel, si consulti al riguardo
  sez.~5.3.1 di \cite{AGL}.}  in modo che sia possibile vederli anche in caso
di blocco totale del sistema (nell'assunzione che la console sia collegata).

In particolare la stampa dei messaggi sulla console è controllata dal
contenuto del file \sysctlfiled{kernel/printk} (o con l'equivalente parametro
di \func{sysctl}) che prevede quattro valori numerici interi: il primo,
\textit{console\_loglevel}, indica la priorità corrente oltre la quale vengono
stampati i messaggi sulla console, il secondo,
\textit{default\_message\_loglevel}, la priorità di default assegnata ai
messaggi che non ne hanno impostata una, il terzo,
\textit{minimum\_console\_level}, il valore minimo che si può assegnare al
primo valore,\footnote{che può essere usato con una delle operazioni di
  gestione che vedremo a breve per ``\textsl{silenziare}'' il kernel.} ed il
quarto, \textit{default\_console\_loglevel}, il valore di
default.\footnote{anch'esso viene usato nelle operazioni di controllo per
  tornare ad un valore predefinito.}

Per la lettura dei messaggi del kernel e la gestione del relativo buffer
circolare esiste una apposita \textit{system call} chiamata anch'essa
 \texttt{syslog}, ma dato il conflitto di nomi questa viene rimappata su
 un'altra funzione di libreria, in particolare nella \acr{glibc} essa viene
 invocata tramite la funzione \funcd{klogctl},\footnote{nella \acr{libc4} e
   nella \acr{libc5} la funzione invece era \code{SYS\_klog}.} il cui prototipo
 è:

\begin{funcproto}{
\fhead{sys/klog.h}
\fdecl{int klogctl(int op, char *buffer, int len)}
\fdesc{Gestisce i messaggi di log del kernel.} 
}

{La funzione ritorna un intero positivo o nullo dipendente dall'operazione
  scelta in caso di successo e $-1$ per un errore, nel qual caso \var{errno}
  assumerà uno dei valori:
  \begin{errlist}
  \item[\errcode{EINVAL}] l'argomento \param{op} non ha un valore valido, o si
    sono specificati valori non validi per gli altri argomenti quando questi
    sono richiesti.
  \item[\errcode{ENOSYS}] il supporto per \texttt{printk} non è stato compilato
    nel kernel.
  \item[\errcode{EPERM}] non si hanno i privilegi richiesti per l'operazione
    richiesta.
  \item[\errcode{ERESTARTSYS}] l'operazione è stata interrotta da un segnale.
   \end{errlist}
}
\end{funcproto}

La funzione prevede che si passi come primo argomento \param{op} un codice
numerico che indica l'operazione richiesta, il secondo argomento deve essere,
per le operazioni che compiono una lettura di dati, l'indirizzo del buffer su
cui copiarli, ed il terzo quanti byte leggere. L'effettivo uso di questi due
argomenti dipende comunque dall'operazione richiesta, ma essi devono essere
comunque specificati, anche quando non servono, nel qual caso verranno
semplicemente ignorati.

\begin{table}[htb]
  \footnotesize
  \centering
  \begin{tabular}[c]{|l|p{10cm}|}
    \hline
    \textbf{Valore}& \textbf{Significato}\\
    \hline
    \hline
    \texttt{0} & apre il log (attualmente non fa niente), \param{buffer}
                 e \param{len} sono ignorati.\\
    \texttt{1} & chiude il log (attualmente non fa niente), \param{buffer}
                 e \param{len} sono ignorati.\\
    \texttt{2} & legge \param{len} byte nel buffer \param{buffer} dal log dei
                 messaggi.\\   
    \texttt{3} & legge \param{len} byte nel buffer \param{buffer} dal buffer
                 circolare dei messaggi.\\
    \texttt{4} & legge \param{len} byte nel buffer \param{buffer} dal buffer
                 circolare dei messaggi e lo svuota.\\
    \texttt{5} & svuota il buffer circolare dei messaggi, \param{buffer}
                 e \param{len} sono ignorati.\\
    \texttt{6} & disabilita la stampa dei messaggi sulla console, \param{buffer}
                 e \param{len} sono ignorati.\\
    \texttt{7} & abilita la stampa dei messaggi sulla console, \param{buffer}
                 e \param{len} sono ignorati.\\
    \texttt{8} & imposta a \param{len} il livello dei messaggi stampati sulla
                 console, \param{buffer} è ignorato.\\ 
    \texttt{9} & ritorna il numero di byte da leggere presenti sul buffer di
                 log, \param{buffer} e \param{len} sono ignorati (dal kernel
                 2.4.10).\\ 
    \texttt{10}& ritorna la dimensione del buffer di log, \param{buffer}
                 e \param{len} sono ignorati (dal kernel 2.6.6).\\
\hline
\end{tabular}
\caption{Valori possibili per l'argomento \param{op} di \func{klogctl}.}
\label{tab:klogctl_operation}
\end{table}

Si sono riportati in tab.~\ref{tab:klogctl_operation} i valori utilizzabili
per \param{op}, con una breve spiegazione della relativa operazione e di come
vengono usati gli altri due argomenti. Come si può notare la funzione è una
sorta di interfaccia comune usata per eseguire operazioni completamente
diverse fra loro.

L'operazione corrispondente al valore 2 di \param{op} consente di leggere un
messaggio dal cosiddetto \textit{log} del kernel. Eseguire questa operazione è
equivalente ad eseguire una lettura dal file
\procfile{/proc/kmsg},\footnote{in realtà è vero l'opposto, è questa funzione
  che viene eseguita quando si legge da questo file.} se non vi sono messaggi
la funzione si blocca in attesa di dati e ritorna soltanto quando questi
diventano disponibili. In tal caso verranno letti ed
estratti\footnote{estratti in quanti i dati del \textit{log} del kernel si
  possono leggere una volta sola, se più processi eseguono l'operazione di
  lettura soltanto uno riceverà i dati, a meno che completata la propria
  operazione di lettura non restino altri messaggi pendenti che a questo punto
  potrebbero essere letti da un altro processo in attesa.} dal log \param{len}
byte che verranno scritti su \param{buffer}; in questo caso il valore di
ritorno di \func{klogctl} corrisponderà al numero di byte ottenuti.

Se invece si usa l'operazione 3 i dati vengono letti dal buffer circolare
usato da \texttt{printk}, che mantiene tutti i messaggi stampati dal kernel
fino al limite delle sue dimensioni, in questo caso i messaggi possono essere
letti più volte. Usando invece l'operazione 4 si richiede di cancellare il
buffer dopo la lettura, che così risulterà vuoto ad una lettura
successiva. Anche con queste operazioni \param{len} indica il numero di byte
da leggere e \param{buffer} l'indirizzo dove leggerli, e la funzione ritorna
il numero di byte effettivamente letti. L'operazione 5 esegue soltanto la
cancellazione del buffer circolare, \param{len} e \param{buffer} sono ignorati
e la funzione ritorna un valore nullo.

Le operazioni corrispondenti ai valori 6, 7 ed 8 consentono di modificare la
priorità oltre la quale i messaggi vengono stampati direttamente sulla
\textit{console} e fanno riferimento ai parametri del kernel gestiti con le
variabili contenute in \sysctlfile{kernel/printk} di cui abbiamo
parlato prima, ed in particolare con 6 si imposta come corrente il valore
minimo della terza variabile (\textit{minimum\_console\_level}), ottenendo
l'effetto di ridurre al minimo i messaggi che arrivano in console, mentre con
7 si ripristina il valore di default.\footnote{secondo la documentazione
  questo sarebbe quello indicato della quarta variabile,
  \textit{default\_console\_loglevel} in genere pari a 7, ma alcune prove con
  il programma \texttt{mydmesg} che si trova nei sorgenti allegati alla guida
  rivelano che l'unico effetto di questa operazione è riportare il valore a
  quello precedente se lo si è ridotto al minimo con l'operazione 6.}  Per
impostare direttamente un valore specifico infine si può usare 8, nel qual
caso il valore numerico del livello da impostare deve essere specificato
con \param{len}, che può assumere solo un valore fra 1 e 8.

Infine le due operazioni 9 e 10 consentono di ottenere rispettivamente il
numero di byte ancora non letti dal log del kernel, e la dimensione totale di
questo. Per entrambe i dati sono restituiti come valore di ritorno, e gli
argomento \param{buffer} e \param{len} sono ignorati.

Si tenga presente che la modifica del livello minimo per cui i messaggi
vengono stampati sulla console (operazioni 6, 7 e 8) e la cancellazione del
buffer circolare di \texttt{printk} (operazioni 4 e 5) sono privilegiate; fino
al kernel 2.6.30 era richiesta la capacità \const{CAP\_SYS\_ADMIN}, a partire
dal 2.6.38 detto privilegio è stato assegnato ad una capacità aggiuntiva,
\const{CAP\_SYSLOG}. Tutto questo è stato fatto per evitare che processi
eseguiti all'interno di un sistema di virtualizzazione ``\textsl{leggera}''
(come i \textit{Linux Container} di LXC) che necessitano di
\const{CAP\_SYS\_ADMIN} per operare all'interno del proprio ambiente
ristretto, potessero anche avere la capacità di influire sui log del kernel
al di fuori di questo.

\itindend{syslog}



\section{L'I/O su terminale}
\label{sec:sess_terminal_io}

Benché come ogni altro dispositivo i terminali siano accessibili come file,
essi hanno assunto storicamente, essendo stati a lungo l'unico modo di
accedere al sistema, una loro rilevanza specifica, che abbiamo già avuto modo
di incontrare nella precedente sezione.

Esamineremo qui le peculiarità dell'I/O eseguito sui terminali, che per la
loro particolare natura presenta delle differenze rispetto ai normali file su
disco e agli altri dispositivi.



\subsection{L'architettura dell'I/O su terminale}
\label{sec:term_io_design}

I terminali sono una classe speciale di dispositivi a caratteri (si ricordi la
classificazione di sez.~\ref{sec:file_file_types}). Un terminale ha infatti
una caratteristica che lo contraddistingue da un qualunque altro dispositivo,
è infatti destinato a gestire l'interazione con un utente e deve perciò essere
in grado di fare da terminale di controllo per una sessione; tutto questo
comporta la presenza di una serie di capacità specifiche.

L'interfaccia per i terminali è una delle più oscure e complesse, essendosi
stratificata dagli inizi dei sistemi Unix fino ad oggi. Questo comporta una
grande quantità di opzioni e controlli relativi ad un insieme di
caratteristiche (come ad esempio la velocità della linea) necessarie per
dispositivi, come i terminali seriali, che al giorno d'oggi sono praticamente
in disuso.

Storicamente i primi terminali erano appunto terminali di telescriventi
(\textit{teletype}), da cui deriva sia il nome dell'interfaccia, \textit{TTY},
che quello dei relativi file di dispositivo, che sono sempre della forma
\texttt{/dev/tty*}.\footnote{ciò vale solo in parte per i terminali virtuali,
  essi infatti hanno due lati, un \textit{master}, che può assumere i nomi
  \file{/dev/pty[p-za-e][0-9a-f]} ed un corrispondente \textit{slave} con nome
  \file{/dev/tty[p-za-e][0-9a-f]}.}  Oggi essi includono le porte seriali, le
console virtuali dello schermo, i terminali virtuali che vengono creati come
canali di comunicazione dal kernel e che di solito vengono associati alle
connessioni di rete (ad esempio per trattare i dati inviati con \cmd{telnet} o
\cmd{ssh}).

L'I/O sui terminali si effettua con le stesse modalità dei file normali: si
apre il relativo file di dispositivo, e si leggono e scrivono i dati con le
usuali funzioni di lettura e scrittura. Così se apriamo una console virtuale
avremo che \func{read} leggerà quanto immesso dalla tastiera, mentre
\func{write} scriverà sullo schermo.  In realtà questo è vero solo a grandi
linee, perché non tiene conto delle caratteristiche specifiche dei terminali;
una delle principali infatti è che essi prevedono due modalità di operazione,
dette rispettivamente ``\textsl{modo canonico}'' e ``\textsl{modo non
  canonico}'', che hanno dei comportamenti nettamente
diversi. \index{modo~canonico}\index{modo~non~canonico}

% TODO: inserire qui il comportamento di read relativo all'errore EIO sulla
% lettura in background???

La modalità preimpostata all'apertura del terminale è quella canonica, in cui
le operazioni di lettura vengono sempre effettuate assemblando i dati in una
linea. Questo significa che eseguendo una \func{read} su un terminale in modo
canonico la funzione si bloccherà, anche se si sono scritti dei caratteri,
fintanto che non si preme il tasto di ritorno a capo: a questo punto la linea
sarà completata e la funzione ritornerà leggendola per intero.

Inoltre in modalità canonica alcuni dei caratteri che si scrivono sul
terminale vengono interpretati direttamente dal kernel per compiere operazioni
(come la generazione dei segnali associati al \textit{job control} illustrata
in sez.~\ref{sec:sig_job_control}), e non saranno mai letti dal
dispositivo. Quella canonica è di norma la modalità in cui opera la shell.

Un terminale in modo non canonico invece non effettua nessun accorpamento dei
dati in linee né li interpreta; esso viene di solito usato da quei programmi,
come ad esempio gli editor, che necessitano di poter leggere un carattere alla
volta e che gestiscono al loro interno l'interpretazione dei caratteri
ricevuti impiegandoli opportunamente come comandi o come dati.

\begin{figure}[!htb]
  \centering \includegraphics[width=12cm]{img/term_struct}
  \caption{Struttura interna generica del kernel per l'accesso ai dati di un
    terminale.}
  \label{fig:term_struct}
\end{figure}

Per capire le caratteristiche dell'I/O sui terminali occorre esaminare le
modalità con cui esso viene effettuato. L'accesso, come per tutti i
dispositivi, viene gestito dal kernel, ma per tutti i terminali viene
utilizzata una architettura generica che si è schematizzata in
fig.~\ref{fig:term_struct}.  

Ad ogni terminale sono sempre associate due code
per gestire l'input e l'output, che ne implementano una bufferizzazione
all'interno del kernel che è completamente indipendente dalla eventuale
ulteriore bufferizzazione fornita dall'interfaccia standard dei file.

La coda di ingresso mantiene i caratteri che sono stati letti dal terminale ma
non ancora letti da un processo, la sua dimensione è definita dal parametro di
sistema \const{MAX\_INPUT} (si veda sez.~\ref{sec:sys_file_limits}), che ne
specifica il limite minimo, in realtà la coda può essere più grande e cambiare
dimensione dinamicamente. 

Se è stato abilitato il controllo di flusso in ingresso il kernel emette i
caratteri di STOP e START per bloccare e sbloccare l'ingresso dei dati;
altrimenti i caratteri immessi oltre le dimensioni massime vengono persi; in
alcuni casi il kernel provvede ad inviare automaticamente un avviso (un
carattere di BELL, che provoca un beep) sull'output quando si eccedono le
dimensioni della coda.  

Se è abilitato la modalità canonica i caratteri in ingresso restano nella coda
fintanto che non viene ricevuto un a capo; un altro parametro del sistema,
\const{MAX\_CANON}, specifica la dimensione massima di una riga in modalità
canonica.

La coda di uscita è analoga a quella di ingresso e contiene i caratteri
scritti dai processi ma non ancora inviati al terminale. Se è abilitato il
controllo di flusso in uscita il kernel risponde ai caratteri di START e STOP
inviati dal terminale. Le dimensioni della coda non sono specificate, ma non
hanno molta importanza, in quanto qualora esse vengano eccedute il driver
provvede automaticamente a bloccare la funzione chiamante.



\subsection{La gestione delle caratteristiche di un terminale}
\label{sec:term_attr}

Data le loro peculiarità, fin dalla realizzazione dei primi sistemi unix-like
si è posto il problema di come gestire le caratteristiche specifiche dei
terminali. Storicamente i vari dialetti di Unix hanno utilizzato delle
funzioni specifiche diverse fra loro, ma con la realizzazione dello standard
POSIX.1-2001 è stata effettuata opportuna unificazione delle funzioni
attinenti i terminali, sintetizzando le differenze fra BSD e System V in una
unica interfaccia, che è quella adottata da Linux.

Molte delle funzioni previste dallo standard POSIX.1-2001 prevedono come
argomento un file descriptor, dato che in origine le relative operazioni
venivano effettuate con delle opportune chiamate a \func{ioctl}.  Ovviamente
dette funzioni potranno essere usate correttamente soltanto con dei file
descriptor che corrispondono ad un terminale, in caso contrario lo standard
richiede che venga restituito un errore di \errcode{ENOTTY}.

Per evitare l'errore, ed anche semplicemente per verificare se un file
descriptor corrisponde ad un terminale, è disponibile la funzione
\funcd{isatty}, il cui prototipo è:

\begin{funcproto}{
\fhead{unistd.h}
\fdecl{int isatty(int fd)}
\fdesc{Controlla se un file è un terminale.} 
}

{La funzione ritorna $1$ se \param{fd} è connesso ad un terminale e $0$
  altrimenti, nel qual caso \var{errno} assumerà uno dei valori:
  \begin{errlist}
  \item[\errcode{EBADF}] \param{fd} non è un file descriptor valido.
  \item[\errcode{EINVAL}] \param{fd} non è associato a un terminale (non
    ottempera a POSIX.1-2001 che richiederebbe \errcode{ENOTTY}).
  \end{errlist}
}
\end{funcproto}

Un'altra funzione per avere informazioni su un terminale è \funcd{ttyname},
che permette di ottenere il nome del file di dispositivo del terminale
associato ad un file descriptor; il suo prototipo è:

\begin{funcproto}{
\fhead{unistd.h}
\fdecl{char *ttyname(int fd)}
\fdesc{Restituisce il nome del terminale associato ad un file descriptor.} 
}

{La funzione ritorna il puntatore alla stringa contenente il nome del
  terminale in caso di successo e \val{NULL} per un errore, nel qual caso
  \var{errno} assumerà uno dei valori:
  \begin{errlist}
  \item[\errcode{EBADF}] \param{fd} non è un file descriptor valido.
  \item[\errcode{ENOTTY}] \param{fd} non è associato a un terminale.
  \end{errlist}
}
\end{funcproto}

La funzione restituisce il puntatore alla stringa contenente il nome del file
di dispositivo del terminale associato a \param{fd}, che però è allocata
staticamente e può essere sovrascritta da successive chiamate. Per questo
della funzione esiste anche una versione rientrante, \funcd{ttyname\_r}, che
non presenta il problema dell'uso di una zona di memoria statica; il suo
prototipo è:

\begin{funcproto}{
\fhead{unistd.h}
\fdecl{int ttyname\_r(int fd, char *buff, size\_t len)}
\fdesc{Restituisce il nome del terminale associato ad un file descriptor.} 
}

{La funzione ritorna $0$ in caso di successo e $-1$ per un errore, nel qual
  caso \var{errno} assumerà uno dei valori: 
  \begin{errlist}
    \item[\errcode{ERANGE}] la lunghezza del buffer \param{len} non è
      sufficiente per contenere la stringa restituita.
  \end{errlist}
  ed inoltre \errval{EBADF} ed \errval{ENOTTY} con lo stesso significato di
  \func{ttyname}.}
\end{funcproto}

La funzione prende due argomenti in più, il puntatore \param{buff} alla zona
di memoria in cui l'utente vuole che il risultato venga scritto, che dovrà
essere allocata in precedenza, e la relativa dimensione, \param{len}. Se la
stringa che deve essere restituita, compreso lo zero di terminazione finale,
eccede questa dimensione si avrà un errore.

Una funzione analoga alle precedenti prevista da POSIX.1, che restituisce il
nome di un file di dispositivo, è \funcd{ctermid}, il cui prototipo è:

\begin{funcproto}{
\fhead{stdio.h}
\fdecl{char *ctermid(char *s)}
\fdesc{Restituisce il nome del terminale di controllo del processo.} 
}

{La funzione ritorna il puntatore alla stringa contenente il \textit{pathname}
  del terminale o \val{NULL} se non riesce ad eseguire l'operazione, non sono
  previsti errori.}
\end{funcproto}

La funzione restituisce un puntatore al \textit{pathname} del file di
dispositivo del terminale di controllo del processo chiamante.  Se si passa
come argomento \val{NULL} la funzione restituisce il puntatore ad una stringa
statica che può essere sovrascritta da chiamate successive, e non è
rientrante. Indicando invece un puntatore ad una zona di memoria già allocata
la stringa sarà scritta su di essa, ma in questo caso il buffer preallocato
deve avere una dimensione di almeno \constd{L\_ctermid}
caratteri.\footnote{\const{L\_ctermid} è una delle varie costanti del sistema,
  non trattata esplicitamente in sez.~\ref{sec:sys_characteristics}, che
  indica la dimensione che deve avere una stringa per poter contenere il nome
  di un terminale.}

Si tenga presente che il \textit{pathname} restituito dalla funzione potrebbe
non identificare univocamente il terminale (ad esempio potrebbe essere
\file{/dev/tty}), inoltre non è detto che il processo possa effettivamente
essere in grado di aprire il terminale.

I vari attributi associati ad un terminale vengono mantenuti per ciascuno di
essi in una struttura \struct{termios} che viene usata dalle varie funzioni
dell'interfaccia. In fig.~\ref{fig:term_termios} si sono riportati tutti i
campi della definizione di questa struttura usata in Linux; di questi solo i
primi cinque sono previsti dallo standard POSIX.1, ma le varie implementazioni
ne aggiungono degli altri per mantenere ulteriori informazioni.\footnote{la
  definizione della struttura si trova in \file{bits/termios.h}, da non
  includere mai direttamente; Linux, seguendo l'esempio di BSD, aggiunge i due
  campi \var{c\_ispeed} e \var{c\_ospeed} per mantenere le velocità delle
  linee seriali, ed un campo ulteriore, \var{c\_line} per indicare la
  disciplina di linea.}

\begin{figure}[!htb] 
  \footnotesize \centering
  \begin{minipage}[c]{0.80\textwidth}
    \includestruct{listati/termios.h}
  \end{minipage} 
  \normalsize 
  \caption{La struttura \structd{termios}, che identifica le proprietà di un
    terminale.}
  \label{fig:term_termios}
\end{figure}

% per le definizioni dei dati vedi anche:
% http://www.lafn.org/~dave/linux/termios.txt 

I primi quattro campi sono quattro flag che controllano il comportamento del
terminale; essi sono realizzati come maschera binaria, pertanto il tipo
\typed{tcflag\_t} è di norma realizzato con un intero senza segno di lunghezza
opportuna. I valori devono essere specificati bit per bit, avendo cura di non
modificare i bit su cui non si interviene.

\begin{table}[b!ht]
  \footnotesize
  \centering
  \begin{tabular}[c]{|l|p{10cm}|}
    \hline
    \textbf{Valore}& \textbf{Significato}\\
    \hline
    \hline
    \constd{IGNBRK} & Ignora le condizioni di BREAK sull'input. Una
                      \textit{condizione di BREAK} è definita nel contesto di
                      una trasmissione seriale asincrona come una sequenza di
                      bit nulli più lunga di un byte.\\
    \constd{BRKINT} & Controlla la reazione ad un BREAK quando
                      \const{IGNBRK} non è impostato. Se \const{BRKINT} è
                      impostato il BREAK causa lo scarico delle code, 
                      e se il terminale è il terminale di controllo per un 
                      gruppo in foreground anche l'invio di \signal{SIGINT} ai
                      processi di quest'ultimo. Se invece \const{BRKINT} non è
                      impostato un BREAK viene letto come un carattere
                      NUL, a meno che non sia impostato \const{PARMRK}
                      nel qual caso viene letto come la sequenza di caratteri
                      \texttt{0xFF 0x00 0x00}.\\
    \constd{IGNPAR} & Ignora gli errori di parità, il carattere viene passato
                      come ricevuto. Ha senso solo se si è impostato 
                      \const{INPCK}.\\
    \constd{PARMRK} & Controlla come vengono riportati gli errori di parità. Ha 
                      senso solo se \const{INPCK} è impostato e \const{IGNPAR}
                      no. Se impostato inserisce una sequenza \texttt{0xFF
                      0x00} prima di ogni carattere che presenta errori di
                      parità, se non impostato un carattere con errori di
                      parità viene letto come uno \texttt{0x00}. Se un
                      carattere ha il valore \texttt{0xFF} e \const{ISTRIP} 
                      non è impostato, per evitare ambiguità esso viene sempre
                      riportato come \texttt{0xFF 0xFF}.\\
    \constd{INPCK}  & Abilita il controllo di parità in ingresso. Se non viene
                      impostato non viene fatto nessun controllo ed i caratteri
                      vengono passati in input direttamente.\\
    \constd{ISTRIP} & Se impostato i caratteri in input sono tagliati a sette
                      bit mettendo a zero il bit più significativo, altrimenti 
                      vengono passati tutti gli otto bit.\\
    \constd{INLCR}  & Se impostato in ingresso il carattere di a capo
                      (\verb|'\n'|) viene automaticamente trasformato in un
                      ritorno carrello (\verb|'\r'|).\\
    \constd{IGNCR}  & Se impostato il carattere di ritorno carrello 
                      (\textit{carriage return}, \verb|'\r'|) viene scartato 
                      dall'input. Può essere utile per i terminali che inviano 
                      entrambi i caratteri di ritorno carrello e a capo 
                      (\textit{newline}, \verb|'\n'|).\\
    \constd{ICRNL}  & Se impostato un carattere di ritorno carrello  
                      (\verb|'\r'|) sul terminale viene automaticamente 
                      trasformato in un a capo (\verb|'\n'|) sulla coda di
                      input.\\
    \constd{IUCLC}  & Se impostato trasforma i caratteri maiuscoli dal
                      terminale in minuscoli sull'ingresso (opzione non 
                      POSIX).\\
    \constd{IXON}   & Se impostato attiva il controllo di flusso in uscita con i
                      caratteri di START e STOP. se si riceve
                      uno STOP l'output viene bloccato, e viene fatto
                      ripartire solo da uno START, e questi due
                      caratteri non vengono passati alla coda di input. Se non
                      impostato i due caratteri sono passati alla coda di input
                      insieme agli altri.\\
    \constd{IXANY}  & Se impostato con il controllo di flusso permette a
                      qualunque carattere di far ripartire l'output bloccato da
                      un carattere di STOP.\\
    \constd{IXOFF}  & Se impostato abilita il controllo di flusso in
                      ingresso. Il computer emette un carattere di STOP per
                      bloccare l'input dal terminale e lo sblocca con il
                      carattere START.\\
    \constd{IMAXBEL}& Se impostato fa suonare il cicalino se si riempie la cosa
                      di ingresso; in Linux non è implementato e il kernel si
                      comporta cose se fosse sempre impostato (è una estensione
                      BSD).\\
    \constd{IUTF8}  & Indica che l'input è in UTF-8, cosa che consente di
                      utilizzare la cancellazione dei caratteri in maniera
                      corretta (dal kernel 2.6.4 e non previsto in POSIX).\\
    \hline
  \end{tabular}
  \caption{Costanti identificative dei vari bit del flag di controllo
    \var{c\_iflag} delle modalità di input di un terminale.}
  \label{tab:sess_termios_iflag}
\end{table}

Il primo flag, mantenuto nel campo \var{c\_iflag}, è detto \textsl{flag di
  input} e controlla le modalità di funzionamento dell'input dei caratteri sul
terminale, come il controllo di parità, il controllo di flusso, la gestione
dei caratteri speciali; un elenco dei vari bit, del loro significato e delle
costanti utilizzate per identificarli è riportato in
tab.~\ref{tab:sess_termios_iflag}.

Si noti come alcuni di questi flag (come quelli per la gestione del flusso)
fanno riferimento a delle caratteristiche che ormai sono completamente
obsolete; la maggior parte inoltre è tipica di terminali seriali, e non ha
alcun effetto su dispositivi diversi come le console virtuali o gli
pseudo-terminali usati nelle connessioni di rete.

Il secondo flag, mantenuto nel campo \var{c\_oflag}, è detto \textsl{flag di
  output} e controlla le modalità di funzionamento dell'output dei caratteri,
come l'impacchettamento dei caratteri sullo schermo, la traslazione degli a
capo, la conversione dei caratteri speciali; un elenco dei vari bit, del loro
significato e delle costanti utilizzate per identificarli è riportato in
tab.~\ref{tab:sess_termios_oflag}, di questi solo \const{OPOST} era previsto
da POSIX.1, buona parte degli altri sono stati aggiunti in POSIX.1-2001,
quelli ancora assenti sono stati indicati esplicitamente.

\begin{table}[!htb]
  \footnotesize
  \centering
  \begin{tabular}[c]{|l|p{10cm}|}
    \hline
    \textbf{Valore}& \textbf{Significato}\\
    \hline
    \hline
    \constd{OPOST} & Se impostato i caratteri vengono convertiti opportunamente
                    (in maniera dipendente dall'implementazione) per la 
                    visualizzazione sul terminale, ad esempio al
                    carattere di a capo (NL) può venire aggiunto un ritorno
                    carrello (CR).\\
    \constd{OLCUC} & Se impostato trasforma i caratteri minuscoli in ingresso 
                    in caratteri maiuscoli sull'uscita (non previsto da
                    POSIX).\\
    \constd{ONLCR} & Se impostato converte automaticamente il carattere di a 
                    capo (NL) in un carattere di ritorno carrello (CR).\\
    \constd{OCRNL} & Se impostato converte automaticamente il carattere di a
                    capo (NL) nella coppia di caratteri ritorno carrello, a 
                    capo (CR-NL).\\
    \constd{ONOCR} & Se impostato converte il carattere di ritorno carrello
                    (CR) nella coppia di caratteri CR-NL.\\
    \constd{ONLRET}& Se impostato rimuove dall'output il carattere di ritorno
                    carrello (CR).\\
    \constd{OFILL} & Se impostato in caso di ritardo sulla linea invia dei
                    caratteri di riempimento invece di attendere.\\
    \constd{OFDEL} & Se impostato il carattere di riempimento è DEL
                    (\texttt{0x3F}), invece che NUL (\texttt{0x00}), (non
                    previsto da POSIX e non implementato su Linux).\\
    \constd{NLDLY} & Maschera per i bit che indicano il ritardo per il
                    carattere di a capo (NL), i valori possibili sono 
                    \val{NL0} o \val{NL1}.\\
    \constd{CRDLY} & Maschera per i bit che indicano il ritardo per il
                    carattere ritorno carrello (CR), i valori possibili sono
                    \val{CR0}, \val{CR1}, \val{CR2} o \val{CR3}.\\
    \constd{TABDLY}& Maschera per i bit che indicano il ritardo per il
                    carattere di tabulazione, i valori possibili sono
                    \val{TAB0}, \val{TAB1}, \val{TAB2} o \val{TAB3}.\\
    \constd{BSDLY} & Maschera per i bit che indicano il ritardo per il
                    carattere di ritorno indietro (\textit{backspace}), i
                    valori possibili sono \val{BS0} o \val{BS1}.\\
    \constd{VTDLY} & Maschera per i bit che indicano il ritardo per il
                    carattere di tabulazione verticale, i valori possibili sono
                    \val{VT0} o \val{VT1}.\\
    \constd{FFDLY} & Maschera per i bit che indicano il ritardo per il
                    carattere di pagina nuova (\textit{form feed}), i valori
                    possibili sono \val{FF0} o \val{FF1}.\\
    \hline
  \end{tabular}
  \caption{Costanti identificative dei vari bit del flag di controllo
    \var{c\_oflag} delle modalità di output di un terminale.}
  \label{tab:sess_termios_oflag}
\end{table}

Si noti come alcuni dei valori riportati in tab.~\ref{tab:sess_termios_oflag}
fanno riferimento a delle maschere di bit; essi infatti vengono utilizzati per
impostare alcuni valori numerici relativi ai ritardi nell'output di alcuni
caratteri: una caratteristica originaria dei primi terminali su telescrivente,
che avevano bisogno di tempistiche diverse per spostare il carrello in
risposta ai caratteri speciali, e che oggi sono completamente in disuso.

Si tenga presente inoltre che nel caso delle maschere il valore da inserire in
\var{c\_oflag} deve essere fornito avendo cura di cancellare prima tutti i bit
della maschera, i valori da immettere infatti (quelli riportati nella
spiegazione corrispondente) sono numerici e non per bit, per cui possono
sovrapporsi fra di loro. Occorrerà perciò utilizzare un codice del tipo:

\includecodesnip{listati/oflag.c}

\noindent che prima cancella i bit della maschera in questione e poi imposta
il valore.


\begin{table}[!htb]
  \footnotesize
  \centering
  \begin{tabular}[c]{|l|p{10cm}|}
    \hline
    \textbf{Valore}& \textbf{Significato}\\
    \hline
    \hline
    \constd{CBAUD}  & Maschera dei bit (4+1) usati per impostare della velocità
                      della linea (il \textit{baud rate}) in ingresso; non è
                      presente in POSIX ed in Linux non è implementato in
                      quanto viene usato un apposito campo di
                      \struct{termios}.\\ 
    \constd{CBAUDEX}& Bit aggiuntivo per l'impostazione della velocità della
                      linea, non è presente in POSIX e per le stesse
                      motivazioni del precedente non è implementato in Linux.\\
    \constd{CSIZE}  & Maschera per i bit usati per specificare la dimensione 
                      del carattere inviato lungo la linea di trasmissione, i
                      valore ne indica la lunghezza (in bit), ed i valori   
                      possibili sono \val{CS5}, \val{CS6}, \val{CS7} e \val{CS8}
                      corrispondenti ad un analogo numero di bit.\\
    \constd{CSTOPB} & Se impostato vengono usati due bit di stop sulla linea
                      seriale, se non impostato ne viene usato soltanto uno.\\
    \constd{CREAD}  & Se è impostato si può leggere l'input del terminale,
                      altrimenti i caratteri in ingresso vengono scartati
                      quando arrivano.\\
    \constd{PARENB} & Se impostato abilita la generazione il controllo di
                      parità. La reazione in caso di errori dipende dai
                      relativi valori per \var{c\_iflag}, riportati in 
                      tab.~\ref{tab:sess_termios_iflag}. Se non è impostato i
                      bit di parità non vengono generati e i caratteri non
                      vengono controllati.\\
    \constd{PARODD} & Ha senso solo se è attivo anche \const{PARENB}. Se 
                      impostato viene usata una parità è dispari, altrimenti 
                      viene usata una parità pari.\\
    \constd{HUPCL}  & Se è impostato viene distaccata la connessione del
                      modem quando l'ultimo dei processi che ha ancora un file
                      aperto sul terminale lo chiude o esce.\\
    \constd{LOBLK}  & Se impostato blocca l'output su un strato di shell non
                      corrente, non è presente in POSIX e non è implementato
                      da Linux.\\
    \constd{CLOCAL} & Se impostato indica che il terminale è connesso in locale
                      e che le linee di controllo del modem devono essere
                      ignorate. Se non impostato effettuando una chiamata ad
                      \func{open} senza aver specificato il flag di
                      \const{O\_NONBLOCK} si bloccherà il processo finché 
                      non si è stabilita una connessione con il modem; inoltre 
                      se viene rilevata una disconnessione viene inviato un
                      segnale di \signal{SIGHUP} al processo di controllo del
                      terminale. La lettura su un terminale sconnesso comporta
                      una condizione di \textit{end of file} e la scrittura un
                      errore di \errcode{EIO}.\\
    \constd{CIBAUD} & Maschera dei bit della velocità della linea in
                      ingresso; analogo a \const{CBAUD}, non è previsto da
                      POSIX e non è implementato in Linux dato che è
                      mantenuto in un apposito campo di \struct{termios}.\\
    \constd{CMSPAR} & imposta un bit di parità costante: se \const{PARODD} è
                      impostato la parità è sempre 1 (\textit{MARK}) se non è
                      impostato la parità è sempre 0 (\textit{SPACE}), non è
                      previsto da POSIX.\\ 
% vedi: http://www.lothosoft.ch/thomas/libmip/markspaceparity.php
    \constd{CRTSCTS}& Abilita il controllo di flusso hardware sulla seriale,
                      attraverso l'utilizzo delle dei due fili di RTS e CTS.\\
    \hline
  \end{tabular}
  \caption{Costanti identificative dei vari bit del flag di controllo
    \var{c\_cflag} delle modalità di controllo di un terminale.}
  \label{tab:sess_termios_cflag}
\end{table}

Il terzo flag, mantenuto nel campo \var{c\_cflag}, è detto \textsl{flag di
  controllo} ed è legato al funzionamento delle linee seriali, permettendo di
impostarne varie caratteristiche, come il numero di bit di stop, le
impostazioni della parità, il funzionamento del controllo di flusso; esso ha
senso solo per i terminali connessi a linee seriali. Un elenco dei vari bit,
del loro significato e delle costanti utilizzate per identificarli è riportato
in tab.~\ref{tab:sess_termios_cflag}.

I valori previsti per questo flag sono molto specifici, e completamente
attinenti al controllo delle modalità operative di un terminale che opera
attraverso una linea seriale; essi pertanto non hanno nessuna rilevanza per i
terminali che usano un'altra interfaccia fisica, come le console virtuali e
gli pseudo-terminali usati dalle connessioni di rete.

Inoltre alcuni valori di questi flag sono previsti solo per quelle
implementazioni (lo standard POSIX non specifica nulla riguardo
l'implementazione, ma solo delle funzioni di lettura e scrittura) che
mantengono le velocità delle linee seriali all'interno dei flag; come
accennato in Linux questo viene fatto (seguendo l'esempio di BSD) attraverso
due campi aggiuntivi, \var{c\_ispeed} e \var{c\_ospeed}, nella struttura
\struct{termios} (mostrati in fig.~\ref{fig:term_termios}).

\begin{table}[b!ht]
  \footnotesize
  \centering
  \begin{tabular}[c]{|l|p{10cm}|}
    \hline
    \textbf{Valore}& \textbf{Significato}\\
    \hline
    \hline
    \constd{ISIG}   & Se impostato abilita il riconoscimento dei caratteri
                      INTR, QUIT, e SUSP generando il relativo segnale.\\
    \constd{ICANON} & Se impostato il terminale opera in modalità canonica,
                      altrimenti opera in modalità non canonica.\\
    \constd{XCASE}  & Se impostato il terminale funziona solo con le
                      maiuscole. L'input è convertito in minuscole tranne per i
                      caratteri preceduti da una ``\texttt{\bslash}''. In output
                      le maiuscole sono precedute da una ``\texttt{\bslash}'' e 
                      le minuscole convertite in maiuscole. Non è presente in
                      POSIX.\\
    \constd{ECHO}   & Se è impostato viene attivato l'eco dei caratteri in
                      input sull'output del terminale.\\
    \constd{ECHOE}  & Se è impostato l'eco mostra la cancellazione di un
                      carattere in input (in reazione al carattere ERASE)
                      cancellando l'ultimo carattere della riga corrente dallo
                      schermo; altrimenti il carattere è rimandato in eco per
                      mostrare quanto accaduto (usato per i terminali con
                      l'uscita su una stampante).\\
    \constd{ECHOK}  & Se impostato abilita il trattamento della visualizzazione
                      del carattere KILL, andando a capo dopo aver visualizzato
                      lo stesso, altrimenti viene solo mostrato il carattere e
                      sta all'utente ricordare che l'input precedente è stato
                      cancellato.\\
    \constd{ECHONL} & Se impostato viene effettuato l'eco di un a
                      capo (\verb|\n|) anche se non è stato impostato
                      \const{ECHO}.\\
    \constd{ECHOCTL}& Se impostato insieme ad \const{ECHO} i caratteri di
                      controllo ASCII (tranne TAB, NL, START, e STOP) sono
                      mostrati nella forma che prepone un ``\texttt{\circonf}'' 
                      alla lettera ottenuta sommando \texttt{0x40} al valore del
                      carattere (di solito questi si possono ottenere anche
                      direttamente premendo il tasto \texttt{ctrl} più la
                      relativa lettera). Non è presente in POSIX.\\
    \constd{ECHOPRT}& Se impostato abilita la visualizzazione del carattere di
                      cancellazione in una modalità adatta ai terminali con
                      l'uscita su stampante; l'invio del carattere di ERASE
                      comporta la stampa di un ``\texttt{|}'' seguito dal
                      carattere cancellato, e così via in caso di successive
                      cancellazioni, quando si riprende ad immettere carattere 
                      normali prima verrà stampata una ``\texttt{/}''. Non è
                      presente in POSIX.\\
    \constd{ECHOKE} & Se impostato abilita il trattamento della visualizzazione
                      del carattere KILL cancellando i caratteri precedenti
                      nella linea secondo le modalità specificate dai valori di
                      \const{ECHOE} e \const{ECHOPRT}. Non è presente in
                      POSIX.\\ 
    \constd{DEFECHO}& Se impostato effettua l'eco solo se c'è un processo in
                      lettura. Non è presente in POSIX e non è supportato da
                      Linux.\\
    \constd{FLUSHO} & Effettua la cancellazione della coda di uscita. Viene
                      attivato dal carattere DISCARD. Non è presente in POSIX e
                      non è supportato da Linux.\\
    \constd{NOFLSH} & Se impostato disabilita lo scarico delle code di ingresso
                      e uscita quando vengono emessi i segnali \signal{SIGINT}, 
                      \signal{SIGQUIT} e \signal{SIGSUSP}.\\
    \constd{TOSTOP} & Se abilitato, con il supporto per il job control presente,
                      genera il segnale \signal{SIGTTOU} per un processo in
                      background che cerca di scrivere sul terminale.\\
    \constd{PENDIN} & Indica che la linea deve essere ristampata, viene
                      attivato dal carattere REPRINT e resta attivo fino alla
                      fine della ristampa. Non è presente in POSIX e
                      non è supportato in Linux.\\
    \constd{IEXTEN} & Abilita alcune estensioni previste dalla
                      implementazione. Deve essere impostato perché caratteri
                      speciali come EOL2, LNEXT, REPRINT e WERASE possano
                      essere interpretati.\\
    \hline
  \end{tabular}
  \caption{Costanti identificative dei vari bit del flag di controllo
    \var{c\_lflag} delle modalità locali di un terminale.}
  \label{tab:sess_termios_lflag}
\end{table}

Il quarto flag, mantenuto nel campo \var{c\_lflag}, è detto \textsl{flag
  locale}, e serve per controllare il funzionamento dell'interfaccia fra il
driver e l'utente, come abilitare l'eco, gestire i caratteri di controllo e
l'emissione dei segnali, impostare modo canonico o non canonico. Un elenco dei
vari bit, del loro significato e delle costanti utilizzate per identificarli è
riportato in tab.~\ref{tab:sess_termios_lflag}. Con i terminali odierni
l'unico flag con cui probabilmente si può avere a che fare è questo, in quanto
è con questo che si impostano le caratteristiche generiche comuni a tutti i
terminali.

Si tenga presente che i flag che riguardano le modalità di eco dei caratteri
(\const{ECHOE}, \const{ECHOPRT}, \const{ECHOK}, \const{ECHOKE},
\const{ECHONL}) controllano solo il comportamento della visualizzazione, il
riconoscimento dei vari caratteri dipende dalla modalità di operazione, ed
avviene solo in modo canonico, pertanto questi flag non hanno significato se
non è impostato \const{ICANON}.

Oltre ai vari flag per gestire le varie caratteristiche dei terminali,
\struct{termios} contiene pure il campo \var{c\_cc} che viene usato per
impostare i caratteri speciali associati alle varie funzioni di controllo. Il
numero di questi caratteri speciali è indicato dalla costante \constd{NCCS},
POSIX ne specifica almeno 11, ma molte implementazioni ne definiscono molti
altri.\footnote{in Linux il valore della costante è 32, anche se i caratteri
  effettivamente definiti sono solo 17.}

\begin{table}[!htb]
  \footnotesize
  \centering
  \begin{tabular}[c]{|l|c|c|p{9cm}|}
    \hline
    \textbf{Indice} & \textbf{Valore}&\textbf{Codice} & \textbf{Funzione}\\
    \hline
    \hline
    \constd{VINTR} &\texttt{0x03}&(\texttt{C-c})& Carattere di interrupt, 
                                                  provoca l'emissione di 
                                                  \signal{SIGINT}.\\
    \constd{VQUIT} &\texttt{0x1C}&(\texttt{C-\bslash})& Carattere di uscita,  
                                                        provoca l'emissione di 
                                                        \signal{SIGQUIT}.\\
    \constd{VERASE}&\texttt{0x7f}&DEL,\texttt{C-?}& Carattere di ERASE, cancella
                                                    l'ultimo carattere
                                                    precedente nella linea.\\
    \constd{VKILL} &\texttt{0x15}&(\texttt{C-u})& Carattere di KILL, cancella
                                                  l'intera riga.\\
    \constd{VEOF}  &\texttt{0x04}&(\texttt{C-d})& Carattere di
                                                  \textit{end-of-file}. Causa
                                                  l'invio del contenuto del
                                                  buffer di ingresso al
                                                  processo in lettura anche se
                                                  non è ancora stato ricevuto
                                                  un a capo. Se è il primo
                                                  carattere immesso comporta il
                                                  ritorno di \func{read} con
                                                  zero caratteri, cioè la
                                                  condizione di
                                                  \textit{end-of-file}.\\
    \constd{VMIN}  &     ---     &  ---   & Numero minimo di caratteri per una 
                                            lettura in modo non canonico.\\
    \constd{VEOL}  &\texttt{0x00}& NUL &    Carattere di fine riga. Agisce come
                                         un a capo, ma non viene scartato ed
                                         è letto come l'ultimo carattere
                                         nella riga.\\
    \constd{VTIME} &     ---     &  ---   & Timeout, in decimi di secondo, per
                                            una lettura in modo non canonico.\\
    \constd{VEOL2} &\texttt{0x00}&   NUL  & Ulteriore carattere di fine
                                            riga. Ha lo stesso effetto di
                                            \const{VEOL} ma può essere un
                                            carattere diverso. \\
    \constd{VSWTC} &\texttt{0x00}&   NUL  & Carattere di switch. Non supportato
                                            in Linux.\\
    \constd{VSTART}&\texttt{0x17}&(\texttt{C-q})& Carattere di START. Riavvia un
                                                 output bloccato da uno STOP.\\
    \constd{VSTOP} &\texttt{0x19}&(\texttt{C-s})& Carattere di STOP. Blocca
                                                 l'output fintanto che non
                                                 viene premuto un carattere di
                                                 START.\\
    \constd{VSUSP} &\texttt{0x1A}&(\texttt{C-z})& Carattere di
                                                 sospensione. Invia il segnale
                                                 \signal{SIGTSTP}.\\
    \constd{VDSUSP}&\texttt{0x19}&(\texttt{C-y})& Carattere di sospensione
                                                 ritardata. Invia il segnale 
                                                 \signal{SIGTSTP} quando il
                                                 carattere viene letto dal
                                                 programma, (non presente in
                                                 POSIX e non riconosciuto in
                                                 Linux).\\ 
    \constd{VLNEXT}&\texttt{0x16}&(\texttt{C-v})& Carattere di escape, serve a
                                                 quotare il carattere
                                                 successivo che non viene
                                                 interpretato ma passato
                                                 direttamente all'output.\\
    \constd{VWERASE}&\texttt{0x17}&(\texttt{C-w})&Cancellazione di una
                                                 parola.\\
    \constd{VREPRINT}&\texttt{0x12}&(\texttt{C-r})& Ristampa i caratteri non
                                                 ancora letti (non presente in
                                                 POSIX).\\
    \constd{VDISCARD}&\texttt{0x0F}&(\texttt{C-o})& Non riconosciuto in Linux.\\
    \constd{VSTATUS} &\texttt{0x13}&(\texttt{C-t})& Non riconosciuto in Linux.\\
    \hline
  \end{tabular}
  \caption{Valori dei caratteri di controllo mantenuti nel campo \var{c\_cc}
    della struttura \struct{termios}.} 
  \label{tab:sess_termios_cc}
\end{table}


A ciascuna di queste funzioni di controllo corrisponde un elemento del vettore
\var{c\_cc} che specifica quale è il carattere speciale associato; per
portabilità invece di essere indicati con la loro posizione numerica nel
vettore, i vari elementi vengono indicizzati attraverso delle opportune
costanti, il cui nome corrisponde all'azione ad essi associata. Un elenco
completo dei caratteri di controllo, con le costanti e delle funzionalità
associate è riportato in tab.~\ref{tab:sess_termios_cc}, usando quelle
definizioni diventa possibile assegnare un nuovo carattere di controllo con un
codice del tipo:
\includecodesnip{listati/value_c_cc.c}

La maggior parte di questi caratteri (tutti tranne \const{VTIME} e
\const{VMIN}) hanno effetto solo quando il terminale viene utilizzato in modo
canonico; per alcuni devono essere soddisfatte ulteriori richieste, ad esempio
\const{VINTR}, \const{VSUSP}, e \const{VQUIT} richiedono sia impostato
\const{ISIG}; \const{VSTART} e \const{VSTOP} richiedono sia impostato
\const{IXON}; \const{VLNEXT}, \const{VWERASE}, \const{VREPRINT} richiedono sia
impostato \const{IEXTEN}.  In ogni caso quando vengono attivati i caratteri
vengono interpretati e non sono passati sulla coda di ingresso.

Per leggere ed scrivere tutte le varie impostazioni dei terminali viste finora
lo standard POSIX prevede due funzioni che utilizzano come argomento un
puntatore ad una struttura \struct{termios} che sarà quella in cui andranno
immagazzinate le impostazioni.  Le funzioni sono \funcd{tcgetattr} e
\funcd{tcsetattr} ed il loro prototipo è:

\begin{funcproto}{
\fhead{unistd.h}
\fhead{termios.h}

\fdecl{int tcgetattr(int fd, struct termios *termios\_p)} 
\fdesc{Legge il valore delle impostazioni di un terminale.} 
\fdecl{int tcsetattr(int fd, int optional\_actions, struct termios *termios\_p)}
\fdesc{Scrive le impostazioni di un terminale.}
}

{Le funzioni ritornano $0$ in caso di successo e $-1$ per un errore, nel qual
  caso \var{errno} assumerà uno dei valori: 
  \begin{errlist}
    \item[\errcode{EINTR}] la funzione è stata interrotta. 
  \end{errlist}
  ed inoltre \errval{EBADF}, \errval{ENOTTY} ed \errval{EINVAL} nel loro
  significato generico.}
\end{funcproto}


Le funzioni operano sul terminale cui fa riferimento il file descriptor
\param{fd} utilizzando la struttura indicata dal puntatore \param{termios\_p}
per lo scambio dei dati. Si tenga presente che le impostazioni sono associate
al terminale e non al file descriptor; questo significa che se si è cambiata
una impostazione un qualunque altro processo che apra lo stesso terminale, od
un qualunque altro file descriptor che vi faccia riferimento, vedrà le nuove
impostazioni pur non avendo nulla a che fare con il file descriptor che si è
usato per effettuare i cambiamenti.

Questo significa che non è possibile usare file descriptor diversi per
utilizzare automaticamente il terminale in modalità diverse, se esiste una
necessità di accesso differenziato di questo tipo occorrerà cambiare
esplicitamente la modalità tutte le volte che si passa da un file descriptor
ad un altro.

La funzione \func{tcgetattr} legge i valori correnti delle impostazioni di un
terminale qualunque nella struttura puntata da \param{termios\_p};
\func{tcsetattr} invece effettua la scrittura delle impostazioni e quando
viene invocata sul proprio terminale di controllo può essere eseguita con
successo solo da un processo in foreground.  Se invocata da un processo in
background infatti tutto il gruppo riceverà un segnale di \signal{SIGTTOU} come
se si fosse tentata una scrittura, a meno che il processo chiamante non abbia
\signal{SIGTTOU} ignorato o bloccato, nel qual caso l'operazione sarà eseguita.

La funzione \func{tcsetattr} prevede tre diverse modalità di funzionamento,
specificabili attraverso l'argomento \param{optional\_actions}, che permette
di stabilire come viene eseguito il cambiamento delle impostazioni del
terminale, i valori possibili sono riportati in
tab.~\ref{tab:sess_tcsetattr_option}; di norma (come fatto per le due funzioni
di esempio) si usa sempre \const{TCSANOW}, le altre opzioni possono essere
utili qualora si cambino i parametri di output.

\begin{table}[htb]
  \footnotesize
  \centering
  \begin{tabular}[c]{|l|p{8cm}|}
    \hline
    \textbf{Valore}& \textbf{Significato}\\
    \hline
    \hline
    \constd{TCSANOW}  & Esegue i cambiamenti in maniera immediata.\\
    \constd{TCSADRAIN}& I cambiamenti vengono eseguiti dopo aver atteso che
                        tutto l'output presente sulle code è stato scritto.\\
    \constd{TCSAFLUSH}& È identico a \const{TCSADRAIN}, ma in più scarta
                        tutti i dati presenti sulla coda di input.\\
    \hline
  \end{tabular}
  \caption{Possibili valori per l'argomento \param{optional\_actions} della
    funzione \func{tcsetattr}.} 
  \label{tab:sess_tcsetattr_option}
\end{table}

Occorre infine tenere presente che \func{tcsetattr} ritorna con successo anche
se soltanto uno dei cambiamenti richiesti è stato eseguito. Pertanto se si
effettuano più cambiamenti è buona norma controllare con una ulteriore
chiamata a \func{tcgetattr} che essi siano stati eseguiti tutti quanti.

\begin{figure}[!htbp]
  \footnotesize \centering
  \begin{minipage}[c]{\codesamplewidth}
    \includecodesample{listati/SetTermAttr.c}
  \end{minipage} 
  \normalsize 
  \caption{Codice della funzione \func{SetTermAttr} che permette di
    impostare uno dei flag di controllo locale del terminale.}
  \label{fig:term_set_attr}
\end{figure}

Come già accennato per i cambiamenti effettuati ai vari flag di controllo
occorre che i valori di ciascun bit siano specificati avendo cura di mantenere
intatti gli altri; per questo motivo in generale si deve prima leggere il
valore corrente delle impostazioni con \func{tcgetattr} per poi modificare i
valori impostati.

In fig.~\ref{fig:term_set_attr} e fig.~\ref{fig:term_unset_attr} si è
riportato rispettivamente il codice delle due funzioni \func{SetTermAttr} e
\func{UnSetTermAttr}, che possono essere usate per impostare o rimuovere, con
le dovute precauzioni, un qualunque bit di \var{c\_lflag}. Il codice completo
di entrambe le funzioni può essere trovato nel file \file{SetTermAttr.c} dei
sorgenti allegati alla guida.

La funzione \func{SetTermAttr} provvede ad impostare il bit specificato
dall'argomento \param{flag}; prima si leggono i valori correnti
(\texttt{\small 8}) con \func{tcgetattr}, uscendo con un messaggio in caso di
errore (\texttt{\small 9--10}), poi si provvede a impostare solo i bit
richiesti (possono essere più di uno) con un OR binario (\texttt{\small 12});
infine si scrive il nuovo valore modificato con \func{tcsetattr}
(\texttt{\small 13}), notificando un eventuale errore (\texttt{\small 14--15})
o uscendo normalmente.

\begin{figure}[!htbp]
  \footnotesize \centering
  \begin{minipage}[c]{\codesamplewidth}
    \includecodesample{listati/UnSetTermAttr.c}
  \end{minipage} 
  \normalsize 
  \caption{Codice della funzione \func{UnSetTermAttr} che permette di
    rimuovere uno dei flag di controllo locale del terminale.} 
  \label{fig:term_unset_attr}
\end{figure}

La seconda funzione, \func{UnSetTermAttr}, è assolutamente identica alla
prima, solo che in questo caso (\texttt{\small 9}) si rimuovono i bit
specificati dall'argomento \param{flag} usando un AND binario del valore
negato.

Al contrario di tutte le altre caratteristiche dei terminali, che possono
essere impostate esplicitamente utilizzando gli opportuni campi di
\struct{termios}, per le velocità della linea (il cosiddetto \textit{baud
  rate}) non è prevista una implementazione standardizzata, per cui anche se
in Linux sono mantenute in due campi dedicati nella struttura, questi non
devono essere acceduti direttamente ma solo attraverso le apposite funzioni di
interfaccia provviste da POSIX.1.

Lo standard prevede due funzioni per scrivere la velocità delle linee seriali,
\funcd{cfsetispeed} per la velocità della linea di ingresso e
\funcd{cfsetospeed} per la velocità della linea di uscita; i loro prototipi
sono:

\begin{funcproto}{
\fhead{unistd.h}
\fhead{termios.h}
\fdecl{int cfsetispeed(struct termios *termios\_p, speed\_t speed)}
\fdesc{Imposta la velocità delle linee seriali in ingresso.}
\fdecl{int cfsetospeed(struct termios *termios\_p, speed\_t speed)} 
\fdesc{Imposta la velocità delle linee seriali in uscita.} 
}

{Le funzioni ritornano $0$ in caso di successo e $-1$ per un errore, che
  avviene solo quando il valore specificato non è valido.}
\end{funcproto}

Si noti che le funzioni si limitano a scrivere opportunamente il valore della
velocità prescelta \param{speed} all'interno della struttura puntata da
\param{termios\_p}; per effettuare l'impostazione effettiva occorrerà poi
chiamare \func{tcsetattr}.

Si tenga presente che per le linee seriali solo alcuni valori di velocità sono
validi; questi possono essere specificati direttamente (la \acr{glibc}
prevede che i valori siano indicati in bit per secondo), ma in generale
altre versioni di libreria possono utilizzare dei valori diversi. Per questo
POSIX.1 prevede una serie di costanti che però servono solo per specificare le
velocità tipiche delle linee seriali:
\begin{verbatim}
     B0       B50      B75      B110     B134     B150     B200     
     B300     B600     B1200    B1800    B2400    B4800    B9600    
     B19200   B38400   B57600   B115200  B230400  B460800
\end{verbatim}

Un terminale può utilizzare solo alcune delle velocità possibili, le funzioni
però non controllano se il valore specificato è valido, dato che non possono
sapere a quale terminale le velocità saranno applicate; sarà l'esecuzione di
\func{tcsetattr} a fallire quando si cercherà di eseguire l'impostazione.

Di norma il valore ha senso solo per i terminali seriali dove indica appunto
la velocità della linea di trasmissione; se questa non corrisponde a quella
del terminale quest'ultimo non potrà funzionare: quando il terminale non è
seriale il valore non influisce sulla velocità di trasmissione dei dati. 

In generale impostare un valore nullo (\val{B0}) sulla linea di output fa sì
che il modem non asserisca più le linee di controllo, interrompendo di fatto
la connessione, qualora invece si utilizzi questo valore per la linea di input
l'effetto sarà quello di rendere la sua velocità identica a quella della linea
di output.

Dato che in genere si imposta sempre la stessa velocità sulle linee di uscita
e di ingresso è supportata anche la funzione \funcd{cfsetspeed}, una
estensione di BSD (la funzione origina da 4.4BSD e richiede sia definita la
macro \macro{\_BSD\_SOURCE}) il cui prototipo è:

\begin{funcproto}{
\fhead{unistd.h} 
\fhead{termios.h}
\fdecl{int cfsetspeed(struct termios *termios\_p, speed\_t speed)}
\fdesc{Imposta la velocità delle linee seriali.} 
}

{La funzione ritorna $0$ in caso di successo e $-1$ per un errore, che avviene
  solo quando il valore specificato non è valido.}
\end{funcproto}

\noindent la funzione è identica alle due precedenti ma imposta la stessa
velocità sia per la linea di ingresso che per quella di uscita.

Analogamente a quanto avviene per l'impostazione, le velocità possono essere
lette da una struttura \struct{termios} utilizzando altre due funzioni,
\funcd{cfgetispeed} e \funcd{cfgetospeed}, i cui prototipi sono:

\begin{funcproto}{
\fhead{unistd.h} 
\fhead{termios.h}
\fdecl{speed\_t cfgetispeed(struct termios *termios\_p)} 
\fdesc{Legge la velocità delle linee seriali in ingresso.}
\fdecl{speed\_t cfgetospeed(struct termios *termios\_p)} 
\fdesc{Legge la velocità delle linee seriali in uscita.}
}

{Le funzioni ritornano la velocità della linea, non sono previste condizioni
  di errore.}
\end{funcproto}

Anche in questo caso le due funzioni estraggono i valori della velocità della
linea da una struttura, il cui indirizzo è specificato dall'argomento
\param{termios\_p} che deve essere stata letta in precedenza con
\func{tcgetattr}.

Infine sempre da BSD è stata ripresa una funzione che consente di impostare il
terminale in una modalità analoga alla cosiddetta modalità ``\textit{raw}'' di
System V, in cui i dati in input vengono resi disponibili un carattere alla
volta, e l'eco e tutte le interpretazioni dei caratteri in entrata e uscita
sono disabilitate. La funzione è \funcd{cfmakeraw} ed il suo prototipo è:

\begin{funcproto}{
\fhead{unistd.h} 
\fhead{termios.h}
\fdecl{void cfmakeraw(struct termios *termios\_p)} 
\fdesc{Imposta il terminale in  modalità ``\textit{raw}''.}
}

{La funzione imposta solo i valori in \param{termios\_p}, e non
    sono previste condizioni di errore.}
\end{funcproto}

Anche in questo caso la funzione si limita a preparare i valori che poi
saranno impostato con una successiva chiamata a \func{tcsetattr}, in sostanza
la funzione è equivalente a:
\includecodesnip{listati/cfmakeraw.c}


\subsection{La gestione della disciplina di linea.}
\label{sec:term_line_discipline}

Come illustrato dalla struttura riportata in fig.~\ref{fig:term_struct} tutti
i terminali hanno un insieme di funzionalità comuni, che prevedono la presenza
di code di ingresso ed uscita; in generale si fa riferimento a queste
funzionalità con il nome di \textsl{disciplina di
  linea}.\index{disciplina~di~linea} Lo standard POSIX prevede alcune funzioni
che permettono di intervenire direttamente sulla gestione della disciplina di
linea e sull'interazione fra i dati in ingresso ed uscita e le relative code.

In generale tutte queste funzioni vengono considerate, dal punto di vista
dell'accesso al terminale, come delle funzioni di scrittura, pertanto se usate
da processi in background sul loro terminale di controllo provocano
l'emissione di \signal{SIGTTOU}, come illustrato in
sez.~\ref{sec:sess_ctrl_term}, con la stessa eccezione, già vista per
\func{tcsetattr}, che quest'ultimo sia bloccato o ignorato dal processo
chiamante.

Una prima funzione, che è efficace solo in caso di terminali seriali
asincroni, e non fa niente per tutti gli altri terminali, è
\funcd{tcsendbreak}; il suo prototipo è:

\begin{funcproto}{
\fhead{unistd.h} 
\fhead{termios.h}
\fdecl{int tcsendbreak(int fd, int duration)}
\fdesc{Genera una condizione di \textit{break}.}  

}

{La funzione ritorna $0$ in caso di successo e $-1$ per un errore, nel qual
  caso \var{errno} assumerà uno dei valori \errval{EBADF} o \errval{ENOTTY}
  nel loro significato generico.}
\end{funcproto}

La funzione invia un flusso di bit nulli, che genera una condizione di
\textit{break}, sul terminale associato a \param{fd}. Un valore nullo
di \param{duration} implica una durata del flusso fra 0.25 e 0.5 secondi, un
valore diverso da zero implica una durata pari a \code{duration*T} dove
\code{T} è un valore compreso fra 0.25 e 0.5 secondi. Lo standard POSIX
specifica il comportamento solo nel caso si sia impostato un valore nullo
per \param{duration}, il comportamento negli altri casi può dipendere
dall'implementazione.

Le altre funzioni previste dallo standard POSIX servono a controllare il
comportamento dell'interazione fra le code associate al terminale e l'utente;
la prima di queste è \funcd{tcdrain}, il cui prototipo è:


\begin{funcproto}{
\fhead{unistd.h} 
\fhead{termios.h}
\fdecl{int tcdrain(int fd)}
\fdesc{Attende lo svuotamento della coda di uscita.} 
}

{La funzione ritorna $0$ in caso di successo e $-1$ per un errore, nel qual
  caso \var{errno} assumerà i valori \errval{EBADF} o \errval{ENOTTY}.}
\end{funcproto}

La funzione blocca il processo fino a che tutto l'output presente sulla coda
di uscita non è stato trasmesso al terminale associato ad \param{fd}. % La
                                % funzione è  un punto di cancellazione per i
                                % programmi multi-thread, in tal caso le
                                % chiamate devono essere protette con dei
                                % gestori di cancellazione. 

Una seconda funzione, \funcd{tcflush}, permette svuotare immediatamente le code
di cancellando tutti i dati presenti al loro interno; il suo prototipo è:

\begin{funcproto}{
\fhead{unistd.h}
\fhead{termios.h}
\fdecl{int tcflush(int fd, int queue)}
\fdesc{Cancella i dati presenti nelle code di ingresso o di uscita.}  }

{La funzione ritorna $0$ in caso di successo e $-1$ per un errore, nel qual
  caso \var{errno} assumerà i valori \errval{EBADF} o \errval{ENOTTY}.}
\end{funcproto}

La funzione agisce sul terminale associato a \param{fd}, l'argomento
\param{queue} permette di specificare su quale coda (ingresso, uscita o
entrambe), operare. Esso può prendere i valori riportati in
tab.~\ref{tab:sess_tcflush_queue}, nel caso si specifichi la coda di ingresso
cancellerà i dati ricevuti ma non ancora letti, nel caso si specifichi la coda
di uscita cancellerà i dati scritti ma non ancora trasmessi.

\begin{table}[htb]
  \footnotesize
  \centering
  \begin{tabular}[c]{|l|l|}
    \hline
    \textbf{Valore}& \textbf{Significato}\\
    \hline
    \hline
    \constd{TCIFLUSH} & Cancella i dati sulla coda di ingresso.\\
    \constd{TCOFLUSH} & Cancella i dati sulla coda di uscita. \\
    \constd{TCIOFLUSH}& Cancella i dati su entrambe le code.\\
    \hline
  \end{tabular}
  \caption{Possibili valori per l'argomento \param{queue} della
    funzione \func{tcflush}.} 
  \label{tab:sess_tcflush_queue}
\end{table}


L'ultima funzione dell'interfaccia che interviene sulla disciplina di linea è
\funcd{tcflow}, che viene usata per sospendere la trasmissione e la ricezione
dei dati sul terminale; il suo prototipo è:

\begin{funcproto}{
\fhead{unistd.h}
\fhead{termios.h}
\fdecl{int tcflow(int fd, int action)}
\fdesc{Sospende e riavvia il flusso dei dati sul terminale.} 
}

{La funzione ritorna $0$ in caso di successo e $-1$ per un errore, nel qual
  caso \var{errno} assumerà i valori \errval{EBADF} o \errval{ENOTTY}.}
\end{funcproto}

La funzione permette di controllare (interrompendo e facendo riprendere) il
flusso dei dati fra il terminale ed il sistema sia in ingresso che in uscita.
Il comportamento della funzione è regolato dall'argomento \param{action}, i
cui possibili valori, e relativa azione eseguita dalla funzione, sono
riportati in tab.~\ref{tab:sess_tcflow_action}.

\begin{table}[htb]
   \footnotesize
  \centering
  \begin{tabular}[c]{|l|p{8cm}|}
    \hline
    \textbf{Valore}& \textbf{Azione}\\
    \hline
    \hline
    \constd{TCOOFF}& Sospende l'output.\\
    \constd{TCOON} & Riprende un output precedentemente sospeso.\\
    \constd{TCIOFF}& Il sistema trasmette un carattere di STOP, che 
                     fa interrompere la trasmissione dei dati dal terminale.\\
    \constd{TCION} & Il sistema trasmette un carattere di START, che 
                     fa riprendere la trasmissione dei dati dal terminale.\\
    \hline
  \end{tabular}
  \caption{Possibili valori per l'argomento \param{action} della
    funzione \func{tcflow}.} 
  \label{tab:sess_tcflow_action}
\end{table}



\subsection{Operare in \textsl{modo non canonico}}
\label{sec:term_non_canonical}

Operare con un terminale in modo canonico è relativamente semplice; basta
eseguire una lettura e la funzione ritornerà quando il terminale avrà
completato una linea di input. Non è detto che la linea sia letta interamente
(si può aver richiesto un numero inferiore di byte) ma in ogni caso nessun
dato verrà perso, e il resto della linea sarà letto alla chiamata successiva.

Inoltre in modo canonico la gestione dei dati in ingresso è di norma eseguita
direttamente dal kernel, che si incarica (a seconda di quanto impostato con le
funzioni viste nei paragrafi precedenti) di cancellare i caratteri, bloccare e
riavviare il flusso dei dati, terminare la linea quando viene ricevuti uno dei
vari caratteri di terminazione (NL, EOL, EOL2, EOF).

In modo non canonico è invece compito del programma gestire tutto quanto, i
caratteri NL, EOL, EOL2, EOF, ERASE, KILL, CR, REPRINT non vengono
interpretati automaticamente ed inoltre, non dividendo più l'input in linee,
il sistema non ha più un limite definito su quando ritornare i dati ad un
processo. Per questo motivo abbiamo visto che in \var{c\_cc} sono previsti due
caratteri speciali, MIN e TIME (specificati dagli indici \const{VMIN} e
\const{VTIME} in \var{c\_cc}) che dicono al sistema di ritornare da una
\func{read} quando è stata letta una determinata quantità di dati o è passato
un certo tempo.

Come accennato nella relativa spiegazione in tab.~\ref{tab:sess_termios_cc},
TIME e MIN non sono in realtà caratteri ma valori numerici. Il comportamento
del sistema per un terminale in modalità non canonica prevede quattro casi
distinti:
\begin{description}
\item[MIN$>0$, TIME$>0$] In questo caso MIN stabilisce il numero minimo di
  caratteri desiderati e TIME un tempo di attesa, in decimi di secondo, fra un
  carattere e l'altro. Una \func{read} ritorna se vengono ricevuti almeno MIN
  caratteri prima della scadenza di TIME (MIN è solo un limite inferiore, se
  la funzione ha richiesto un numero maggiore di caratteri ne possono essere
  restituiti di più); se invece TIME scade vengono restituiti i byte ricevuti
  fino ad allora (un carattere viene sempre letto, dato che il timer inizia a
  scorrere solo dopo la ricezione del primo carattere).
\item[MIN$>0$, TIME$=0$] Una \func{read} ritorna solo dopo che sono stati
  ricevuti almeno MIN caratteri. Questo significa che una \func{read} può
  bloccarsi indefinitamente. 
\item[MIN$=0$, TIME$>0$] In questo caso TIME indica un tempo di attesa dalla
  chiamata di \func{read}, la funzione ritorna non appena viene ricevuto un
  carattere o scade il tempo. Si noti che è possibile che \func{read} ritorni
  con un valore nullo.
\item[MIN$=0$, TIME$=0$] In questo caso una \func{read} ritorna immediatamente
  restituendo tutti i caratteri ricevuti. Anche in questo caso può ritornare
  con un valore nullo.
\end{description}



\section{La gestione dei terminali virtuali}
\label{sec:sess_virtual_terminal}

% 
% TODO terminali virtuali 
% Qui c'è da mettere tutta la parte sui terminali virtuali, e la gestione
% degli stessi
%

Da fare.

\subsection{I terminali virtuali}
\label{sec:sess_pty}

Qui vanno spiegati i terminali virtuali, \file{/dev/pty} e compagnia.
% vedi man pts
% vedi 


\subsection{Allocazione dei terminali virtuali}
\label{sec:sess_openpty}

Qui vanno le cose su \func{openpty} e compagnia.

% TODO le ioctl dei terminali (man tty_ioctl)
% e http://www.net-security.org/article.php?id=83
% TODO trattare \func{posix\_openpt}
% vedi http://lwn.net/Articles/688809/,
% http://man7.org/linux/man-pages/man3/ptsname.3.html


% TODO materiale sulle seriali
% vedi http://www.easysw.com/~mike/serial/serial.html
% TODO materiale generico sul layer TTY
% vedi http://www.linusakesson.net/programming/tty/index.php

% LocalWords:  kernel multitasking job control BSD POSIX shell sez group
% LocalWords:  foreground process bg fg waitpid WUNTRACED pgrp session sched
% LocalWords:  struct pgid sid pid ps getpgid getpgrp unistd void errno int
% LocalWords:  ESRCH getsid glibc system call XOPEN SOURCE EPERM setpgrp EACCES
% LocalWords:  setpgid exec EINVAL did fork race condition setsid tty ioctl
% LocalWords:  NOCTTY TIOCSCTTY error tcsetpgrp termios fd pgrpid descriptor VT
% LocalWords:  ENOTTY ENOSYS EBADF SIGTTIN SIGTTOU EIO tcgetpgrp crypt SIGTSTP
% LocalWords:  SIGINT SIGQUIT SIGTERM SIGHUP hungup kill orphaned SIGCONT exit
% LocalWords:  init Slackware run level inittab fig device getty TERM at execle
% LocalWords:  getpwnam chdir home chown chmod setuid setgid initgroups SIGCHLD
% LocalWords:  daemon like daemons NdT Stevens Programming FAQ filesystem umask
% LocalWords:  noclose syslog syslogd socket UDP klogd printk printf facility
% LocalWords:  priority log openlog const char ident option argv tab AUTH CRON
% LocalWords:  AUTHPRIV cron FTP KERN LOCAL LPR NEWS news USENET UUCP CONS CRIT
% LocalWords:  NDELAY NOWAIT ODELAY PERROR stderr format strerror EMERG ALERT
% LocalWords:  ERR WARNING NOTICE INFO DEBUG debug setlogmask mask UPTO za ssh
% LocalWords:  teletype telnet read write BELL beep CANON isatty desc ttyname
% LocalWords:  NULL ctermid stdio pathname buff size len ERANGE bits ispeed xFF
% LocalWords:  ospeed line tcflag INPCK IGNPAR PARMRK ISTRIP IGNBRK BREAK NUL
% LocalWords:  BRKINT IGNCR carriage return newline ICRNL INLCR IUCLC IXON NL
% LocalWords:  IXANY IXOFF IMAXBEL iflag OPOST CR OCRNL OLCUC ONLCR ONOCR OFILL
% LocalWords:  ONLRET OFDEL NLDLY CRDLY TABDLY BSDLY backspace BS VTDLY FFDLY
% LocalWords:  form feed FF oflag CLOCAL of HUPCL CREAD CSTOPB PARENB TIOCGPGRP
% LocalWords:  PARODD CSIZE CS CBAUD CBAUDEX CIBAUD CRTSCTS RTS CTS cflag ECHO
% LocalWords:  ICANON ECHOE ERASE ECHOPRT ECHOK ECHOKE ECHONL ECHOCTL ctrl ISIG
% LocalWords:  INTR QUIT SUSP IEXTEN EOL LNEXT REPRINT WERASE NOFLSH and TOSTOP
% LocalWords:  SIGSUSP XCASE DEFECHO FLUSHO DISCARD PENDIN lflag NCCS VINTR EOF
% LocalWords:  interrupt VQUIT VERASE VKILL VEOF VTIME VMIN VSWTC switch VSTART
% LocalWords:  VSTOP VSUSP VEOL VREPRINT VDISCARD VWERASE VLNEXT escape actions
% LocalWords:  tcgetattr tcsetattr EINTR TCSANOW TCSADRAIN TCSAFLUSH speed MIN
% LocalWords:  SetTermAttr UnSetTermAttr cfsetispeed cfsetospeed cfgetispeed ng
% LocalWords:  cfgetospeed quest'ultime tcsendbreak duration break tcdrain list
% LocalWords:  tcflush queue TCIFLUSH TCOFLUSH TCIOFLUSH tcflow action TCOOFF
% LocalWords:  TCOON TCIOFF TCION timer openpty Window nochdir embedded router
% LocalWords:  access point upstart systemd rsyslog vsyslog variadic src linux
% LocalWords:  closelog dmesg sysctl klogctl sys ERESTARTSYS ConsoleKit to CoPy
% LocalWords:  loglevel message libc klog mydmesg CAP ADMIN LXC pipelining UID
% LocalWords:  TIOCSPGRP GID IUTF UTF LOBLK NONBLOCK CMSPAR MARK VDSUSP VSTATUS
% LocalWords:  cfsetspeed raw cfmakeraw

%%% Local Variables: 
%%% mode: latex
%%% TeX-master: "gapil"
%%% End: 
